\documentclass[11pt]{article}
\usepackage[margin=1in]{geometry}
\usepackage{amsmath,amssymb,amsthm}
\usepackage{enumitem}
\usepackage[colorlinks=true,linkcolor=blue,citecolor=blue,urlcolor=blue]{hyperref}

\theoremstyle{plain}
\newtheorem{theorem}{Theorem}[section]
\newtheorem{lemma}[theorem]{Lemma}
\newtheorem{corollary}[theorem]{Corollary}
\newtheorem{proposition}[theorem]{Proposition}

\theoremstyle{definition}
\newtheorem{definition}[theorem]{Definition}
\newtheorem{example}[theorem]{Example}
\newtheorem{problem}[theorem]{Practice Problem}
\newtheorem{remark}[theorem]{Remark}

\newenvironment{solution}{\begin{proof}[Solution]}{\end{proof}}

\newcommand{\ZZ}{\mathbb{Z}}
\newcommand{\QQ}{\mathbb{Q}}
\newcommand{\RR}{\mathbb{R}}
\newcommand{\CC}{\mathbb{C}}
\newcommand{\NN}{\mathbb{N}}
\newcommand{\FF}{\mathbb{F}}
\newcommand{\ord}{\operatorname{ord}}
\newcommand{\pow}{\operatorname{pow}}
\newcommand{\dang}{\measuredangle}
\DeclareMathOperator{\lcm}{lcm}

\title{Lecture Notes for the IMO Syllabus\\
\large A Thorough Course for Olympiad Trainees}
\author{Compiled Lecture Notes}
\date{\today}

\begin{document}
\maketitle
\tableofcontents
\newpage

%======================================================================
\part{Algebra}
%======================================================================
\section{Algebraic Manipulation and Notation}

\subsection{Factoring and Expanding}
Fluency with identities is the bedrock of olympiad algebra. Memorize and,
more importantly, learn to \emph{recognize} the following:
\begin{align*}
a^2-b^2 &= (a-b)(a+b),\\
a^3 \pm b^3 &= (a\pm b)(a^2 \mp ab + b^2),\\
a^n - b^n &= (a-b)(a^{n-1}+a^{n-2}b+\dots+b^{n-1}),\\
a^n + b^n &= (a+b)(a^{n-1}-a^{n-2}b+\dots+b^{n-1}) \quad (n \text{ odd}),\\
(a+b+c)^2 &= a^2+b^2+c^2+2(ab+bc+ca),\\
(a+b+c)^3 &= a^3+b^3+c^3+3(a+b)(b+c)(c+a),\\
a^3+b^3+c^3-3abc &= (a+b+c)(a^2+b^2+c^2-ab-bc-ca)\\
&= \frac{1}{2}(a+b+c)\bigl((a-b)^2+(b-c)^2+(c-a)^2\bigr).
\end{align*}
Two named factorizations appear constantly:
\begin{itemize}
\item \textbf{Sophie Germain's Identity:} $a^4 + 4b^4 = (a^2+2b^2-2ab)(a^2+2b^2+2ab)$.
\item \textbf{Simon's Favorite Factoring Trick (SFFT):} $xy + ax + by + ab = (x+b)(y+a)$.
\end{itemize}

\begin{example}
Show that $n^4 + 4$ is composite for every integer $n \ge 2$.
\end{example}
\begin{solution}
By Sophie Germain with $a=n$, $b=1$:
$n^4+4 = (n^2-2n+2)(n^2+2n+2)$.
For $n\ge 2$ both factors satisfy $n^2-2n+2 = (n-1)^2+1 \ge 2$, so the number is composite.
\end{solution}

\subsection{Sum and Product Notation, Telescoping, and Determinants}
For a sequence $a_1,\dots,a_n$ we write
\[
\sum_{i=1}^n a_i = a_1 + a_2 + \dots + a_n, \qquad
\prod_{i=1}^n a_i = a_1 a_2 \cdots a_n.
\]
Key manipulations include swapping the order of double sums
$\sum_i \sum_j a_{ij} = \sum_j \sum_i a_{ij}$,
index shifts, \textbf{telescoping}:
\[
\sum_{k=1}^n \bigl(f(k+1)-f(k)\bigr) = f(n+1)-f(1), \qquad
\prod_{k=1}^n \frac{f(k+1)}{f(k)} = \frac{f(n+1)}{f(1)},
\]
and \textbf{Abel summation (summation by parts)}: with partial sums $A_k = \sum_{j=1}^k a_j$,
\[
\sum_{k=1}^n a_k b_k = A_n b_n + \sum_{k=1}^{n-1} A_k (b_k - b_{k+1}).
\]

\begin{theorem}[Vandermonde Determinant]
For any $x_1, \dots, x_n$,
\[
\det\begin{pmatrix}
1 & x_1 & x_1^2 & \cdots & x_1^{n-1}\\
1 & x_2 & x_2^2 & \cdots & x_2^{n-1}\\
\vdots & \vdots & \vdots & \ddots & \vdots\\
1 & x_n & x_n^2 & \cdots & x_n^{n-1}
\end{pmatrix} = \prod_{1 \le i < j \le n} (x_j - x_i).
\]
\end{theorem}
\begin{proof}
The determinant is a polynomial of total degree $\binom{n}{2}$ in $x_1,\dots,x_n$ that vanishes whenever $x_i = x_j$ (since two rows become identical). Hence it is divisible by $\prod_{i<j}(x_j - x_i)$. Comparing the coefficient of the main diagonal monomial $x_2^1 x_3^2 \cdots x_n^{n-1}$ on both sides shows the constant factor is $1$.
\end{proof}

\begin{example}
Compute $\displaystyle\sum_{k=1}^{n} \frac{1}{k(k+1)}$.
\end{example}
\begin{solution}
Since $\frac{1}{k(k+1)} = \frac1k - \frac1{k+1}$, the sum telescopes to
$1 - \frac{1}{n+1} = \frac{n}{n+1}$.
\end{solution}

\begin{problem}
Prove that $\displaystyle\prod_{k=2}^{n} \left(1 - \frac{1}{k^2}\right) = \frac{n+1}{2n}$.
\end{problem}
\begin{solution}
Write $1-\frac1{k^2} = \frac{(k-1)(k+1)}{k^2}$. Then
\[
\prod_{k=2}^n \frac{(k-1)(k+1)}{k^2}
= \frac{\prod_{k=2}^n (k-1) \cdot \prod_{k=2}^n (k+1)}{\left(\prod_{k=2}^n k\right)^2}
= \frac{(n-1)! \cdot \frac{(n+1)!}{2}}{(n!)^2}
= \frac{n+1}{2n}. \qedhere
\]
\end{solution}

\section{Functions: Injectivity, Surjectivity, and Functional Equations}

\begin{definition}
A \textbf{function} $f\colon A \to B$ assigns to each $a \in A$ exactly one element $f(a)\in B$.
\begin{itemize}
\item $f$ is \textbf{injective} (one-to-one) if $f(x)=f(y) \implies x=y$.
\item $f$ is \textbf{surjective} (onto) if for every $b\in B$ there is $a\in A$ with $f(a)=b$.
\item $f$ is \textbf{bijective} if it is both; then an inverse $f^{-1}$ exists.
\item $f\colon A \to A$ is an \textbf{involution} if $f(f(x)) = x$ for all $x \in A$. Every involution is automatically bijective.
\end{itemize}
\end{definition}

\subsection{Cauchy's Functional Equation and Classical Variants}

\begin{theorem}[Cauchy]\label{thm:cauchy}
Suppose $f\colon \RR \to \RR$ satisfies $f(x+y)=f(x)+f(y)$ for all $x,y$.
Then $f(q x) = q f(x)$ for all $q \in \QQ$; in particular $f(q) = cq$ on $\QQ$
where $c=f(1)$. If in addition $f$ is continuous, or monotonic, or bounded
on some interval, or nonnegative for $x \ge 0$, then $f(x)=cx$ for all real $x$.
\end{theorem}

\begin{proof}
Setting $x=y=0$ gives $f(0)=0$; setting $y=-x$ gives $f(-x)=-f(x)$.
By induction $f(nx)=nf(x)$ for $n\in\NN$, hence for all $n\in\ZZ$.
For $q=m/n$, $n f(qx) = f(nq x) = f(mx) = m f(x)$, so $f(qx)=\frac{m}{n}f(x)$.
Thus $f(q)=cq$ on $\QQ$ with $c=f(1)$.

Now assume $f$ is nondecreasing (if $f(x) \ge 0$ for $x \ge 0$, then $x \ge y \implies f(x)-f(y) = f(x-y) \ge 0$, reducing to this case).
Fix $x\in\RR$ and take rationals $p_n \uparrow x$, $q_n \downarrow x$.
Then $c\,p_n = f(p_n)\le f(x) \le f(q_n) = c\,q_n$. Letting $n\to\infty$
forces $f(x)=cx$. If instead $f$ is bounded on an interval $[a,b]$, say
$|f|\le M$ there, then $g(x) = f(x)-cx$ is additive, vanishes on $\QQ$, and
is bounded on $[a,b]$; additivity and $g(\QQ)=0$ imply $g$ is bounded on all
of $\RR$ (translate by rationals). A nonzero value $g(x_0)\ne 0$ would give
$g(nx_0)=ng(x_0)$ unbounded, contradiction. Hence $g\equiv 0$.
\end{proof}

\begin{remark}[Classical Variants and Hamel Bases]
Under regularity hypotheses (continuity, monotonicity, or positivity):
\begin{enumerate}[label=(\arabic*)]
\item \textbf{Exponential Cauchy:} $f(x+y)=f(x)f(y) \implies f\equiv 0$ or $f(x)=a^x$ (take $\ln f(x)$ when $f>0$).
\item \textbf{Logarithmic Cauchy:} $f(xy)=f(x)+f(y)$ ($x,y>0$) $\implies f(x)=c\ln x$ (set $g(t)=f(e^t)$).
\item \textbf{Multiplicative Cauchy:} $f(xy)=f(x)f(y)$ ($x,y>0$) $\implies f\equiv 0$ or $f(x)=x^c$.
\item \textbf{Jensen's Equation:} $f\bigl(\frac{x+y}{2}\bigr) = \frac{f(x)+f(y)}{2} \implies g(x) := f(x)-f(0)$ is additive.
\end{enumerate}
Without regularity, pathological solutions exist via Hamel bases of $\RR$ over $\QQ$. On olympiads, one derives regularity from the problem (e.g., $f(x^2) = f(x)^2 \ge 0$ gives $f \ge 0$ on positives, hence monotone).
\end{remark}

\subsection{The Olympiad Functional Equation Playbook}
Standard moves in functional equations:
\begin{itemize}
\item \textbf{Substitute special values:} $x=y=0$, $y=0$, $y=-x$, or values that make arguments equal.
\item \textbf{Prove injectivity or surjectivity:} If $f(\text{expr}) = x$, then $f$ is surjective; if $f(x)$ is isolated on one side, fixing the other variables proves injectivity, allowing you to cancel outer $f$'s ($f(A)=f(B) \implies A=B$).
\item \textbf{Exploit symmetry and associativity:} Swap $x \leftrightarrow y$ or expand $f((x+1)^2)$ in two ways.
\end{itemize}

\begin{problem}
Find all $f\colon\RR\to\RR$ with $f(x^2 - y^2) = x f(x) - y f(y)$ for all $x,y$.
\end{problem}
\begin{solution}
Put $x=y$: $f(0)=0$. Put $y=0$: $f(x^2)=x f(x)$; replacing
$x \to -x$ gives $x f(x) = -x f(-x)$, so $f$ is odd. The equation becomes
$f(x^2-y^2)=f(x^2)-f(y^2)$, i.e.\ $f(a-b)=f(a)-f(b)$ for $a,b \ge 0$, which
with oddness gives full additivity: $f(u+v)=f(u)+f(v)$ for all reals.
Also $f(x^2)=xf(x)$; expand $f((x+1)^2) = (x+1)f(x+1)$ using additivity:
\[
f(x^2) + 2f(x) + f(1) = (x+1)\bigl(f(x)+f(1)\bigr)
= x f(x) + f(x) + x f(1) + f(1).
\]
Since $f(x^2)=xf(x)$, this simplifies to $f(x) = x f(1)$. So $f(x)=cx$, and all
such functions check.
\end{solution}

\section{Polynomials}

\subsection{One Variable: Division, Vieta, and Arithmetic Properties}
A polynomial of degree $n$ is $P(x)=a_nx^n+\dots+a_0$ with $a_n\ne 0$.

\begin{theorem}[Division Algorithm]
For polynomials $P, D$ with $D \ne 0$ over a field, there exist unique $Q,R$
with $P = QD + R$ and $\deg R < \deg D$ (or $R=0$).
\end{theorem}

\begin{corollary}[Factor Theorem]
$P(r)=0 \iff (x-r) \mid P(x)$. Consequently a nonzero polynomial of degree $n$
has at most $n$ roots; two degree-$\le n$ polynomials agreeing at $n+1$ points
are identical.
\end{corollary}

\begin{proof}
Divide $P$ by $x - r$: $P(x)=(x-r)Q(x)+P(r)$. The rest follows by induction on the
number of roots.
\end{proof}

\begin{theorem}[Fundamental Theorem of Algebra]
Every nonconstant polynomial with complex coefficients has a complex root;
hence $P$ of degree $n$ factors as $a_n\prod_{i=1}^n (x-r_i)$ over $\CC$.
\end{theorem}

\begin{proof}[Proof sketch]
$|P(z)|\to\infty$ as $|z|\to\infty$, so $|P|$ attains a minimum at some $z_0$.
If $P(z_0)\ne 0$, write $P(z_0+w) = P(z_0) + c w^k + O(w^{k+1})$ with $c\ne0$ and
choose the direction of $w$ so that $cw^k$ points opposite to $P(z_0)$; for small
$|w|$ this strictly decreases $|P|$, a contradiction (d'Alembert's argument).
\end{proof}

\begin{theorem}[Vieta's Relations]
If $P(x)=a_n \prod_{i=1}^{n}(x - r_i) = \sum_{k=0}^n a_k x^k$, then for each $1 \le k \le n$,
\[
e_k(r_1,\dots,r_n) := \sum_{i_1 < \dots < i_k} r_{i_1}\cdots r_{i_k}
= (-1)^k \frac{a_{n-k}}{a_n}.
\]
\end{theorem}
\begin{proof}
Expand the product $\prod (x - r_i)$: the coefficient of $x^{n-k}$ collects one
factor $-r_i$ from exactly $k$ of the brackets, giving $(-1)^k e_k$. Compare
coefficients.
\end{proof}

\begin{proposition}\label{prop:ab}
If $P \in \ZZ[x]$ and $a,b \in \ZZ$, then $(a-b) \mid P(a)-P(b)$.
\end{proposition}
\begin{proof}
$P(a)-P(b) = \sum_k c_k (a^k - b^k)$ and $(a-b)\mid a^k-b^k$ for each $k$.
\end{proof}

\begin{theorem}[Rational Root Theorem]
If $p/q$ (in lowest terms) is a root of $a_nx^n+\dots+a_0 \in \ZZ[x]$, then
$p \mid a_0$ and $q \mid a_n$.
\end{theorem}

\begin{example}
Let $r,s,t$ be the roots of $x^3 - 2x^2 + 3x - 4$. Compute $r^2+s^2+t^2$.
\end{example}
\begin{solution}
By Vieta, $r+s+t = 2$, $rs+st+tr = 3$. Then
$r^2+s^2+t^2 = (r+s+t)^2 - 2(rs+st+tr) = 4 - 6 = -2$.
\end{solution}

\begin{problem}
Let $P\in\ZZ[x]$ and suppose $P(a)=P(b)=P(c)=2$ for distinct integers $a,b,c$.
Prove $P(x)=3$ has no integer solution.
\end{problem}
\begin{solution}
Write $P(x)-2 = (x-a)(x-b)(x-c)Q(x)$ with $Q\in\ZZ[x]$ (divide successively; integer
coefficients persist since leading divisions are by monic factors). If
$P(d)=3$ for $d \in \ZZ$, then $(d-a)(d-b)(d-c)Q(d)=1$, so each of
$d-a,\,d-b,\,d-c \in \{1,-1\}$. By pigeonhole two of them are equal, contradicting
$a,b,c$ distinct.
\end{solution}

\subsection{Lagrange Interpolation and Irreducibility Criteria}

\begin{theorem}[Lagrange Interpolation]
Given $n+1$ distinct points $x_0, x_1, \dots, x_n$ and values $y_0, y_1, \dots, y_n$ over a field, there exists a unique polynomial $P(x)$ of degree $\le n$ such that $P(x_i) = y_i$ for all $i$:
\[
P(x) = \sum_{i=0}^n y_i \prod_{j \ne i} \frac{x - x_j}{x_i - x_j}.
\]
In particular, the coefficient of $x^n$ in $P(x)$ is $\sum_{i=0}^n \frac{y_i}{\prod_{j \ne i}(x_i - x_j)}$.
\end{theorem}
\begin{proof}
Each basis polynomial $L_i(x) = \prod_{j \ne i} \frac{x-x_j}{x_i-x_j}$ has degree $n$ and satisfies $L_i(x_k) = \delta_{ik}$. Thus $P(x) = \sum_{i=0}^n y_i L_i(x)$ has degree $\le n$ and satisfies $P(x_k) = y_k$. Uniqueness follows from the Factor Theorem applied to the difference of two such polynomials.
\end{proof}

\begin{theorem}[Gauss's Lemma, Eisenstein, Perron, and Cohn]
Let $P(x) = a_n x^n + \dots + a_0 \in \ZZ[x]$.
\begin{enumerate}[label=(\alph*)]
\item \textbf{(Gauss's Lemma)} $P$ is irreducible over $\QQ$ iff it cannot be factored into two nonconstant polynomials in $\ZZ[x]$.
\item \textbf{(Eisenstein's Criterion)} If there exists a prime $p$ such that $p \nmid a_n$, $p \mid a_i$ for all $0 \le i < n$, and $p^2 \nmid a_0$, then $P(x)$ is irreducible over $\QQ$.
\item \textbf{(Perron's Criterion)} If $a_n = 1$, $a_0 \ne 0$, and $|a_{n-1}| > 1 + |a_{n-2}| + \dots + |a_0|$, then $P(x)$ is irreducible over $\QQ$.
\item \textbf{(Cohn's Criterion)} If a prime $p$ is written in base $b \ge 2$ as $p = \sum_{k=0}^n d_k b^k$ ($0 \le d_k < b$), then $P(x) = \sum_{k=0}^n d_k x^k$ is irreducible over $\QQ$.
\end{enumerate}
\end{theorem}
\begin{proof}[Proof of (b)]
Suppose $P(x) = G(x)H(x)$ in $\ZZ[x]$ with $\deg G, \deg H \ge 1$. Modulo $p$, $\overline{P}(x) = \overline{a_n} x^n = \overline{G}(x)\overline{H}(x)$ in $\FF_p[x]$. By unique factorization in $\FF_p[x]$, both $\overline{G}(x)$ and $\overline{H}(x)$ are monomials in $x$, so their constant terms $g_0, h_0$ are divisible by $p$. Then $p^2 \mid g_0 h_0 = a_0$, contradicting $p^2 \nmid a_0$.
\end{proof}

\begin{problem}
Let $P(x)$ be a polynomial of degree $n$ such that $P(k) = \frac{k}{k+1}$ for $k = 0, 1, \dots, n$. Find $P(n+1)$.
\end{problem}
\begin{solution}
Consider $Q(x) = (x+1)P(x) - x$. Then $\deg Q = n+1$ and $Q(k) = 0$ for $k = 0, 1, \dots, n$. Thus
\[
(x+1)P(x) - x = c\, x(x-1)(x-2)\cdots (x-n).
\]
Plugging in $x = -1$ gives $1 = c(-1)^{n+1}(n+1)!$, so $c = \frac{(-1)^{n+1}}{(n+1)!}$. Evaluating at $x = n+1$ yields
\[
(n+2)P(n+1) - (n+1) = (-1)^{n+1} \implies P(n+1) = \frac{n+1 + (-1)^{n+1}}{n+2}. \qedhere
\]
\end{solution}

\subsection{Symmetric Polynomials and Chebyshev Polynomials}
A polynomial $P(x_1,\dots,x_n)$ is \textbf{symmetric} if it is invariant under
every permutation of the variables, and \textbf{homogeneous of degree $d$} if
every monomial has total degree $d$.

\begin{theorem}[Fundamental Theorem of Symmetric Polynomials]
Every symmetric polynomial in $x_1,\dots,x_n$ can be written uniquely as a
polynomial in the elementary symmetric polynomials $e_1,\dots,e_n$.
\end{theorem}
\begin{proof}[Proof sketch]
Order monomials lexicographically. Given a symmetric $P$ with leading monomial
$x_1^{a_1}\cdots x_n^{a_n}$ (necessarily $a_1\ge\dots\ge a_n$), subtract
$c\, e_1^{a_1-a_2} e_2^{a_2-a_3}\cdots e_n^{a_n}$, which has the same leading
monomial. The leading term strictly decreases; iterate until zero.
\end{proof}

\textbf{Newton's Identities.} With power sums $p_k = \sum x_i^k$ and $e_j = 0$ for $j > n$:
\[
p_k - e_1 p_{k-1} + e_2 p_{k-2} - \dots + (-1)^{k-1}e_{k-1}p_1 + (-1)^{k} k\, e_k = 0 \quad (k \ge 1).
\]

\begin{definition}[Chebyshev Polynomials]
The \textbf{Chebyshev polynomials} $T_n(x)$ are defined by $T_0(x) = 1$, $T_1(x) = x$, and $T_{n+1}(x) = 2xT_n(x) - T_{n-1}(x)$. By induction and $\cos(n+1)\theta + \cos(n-1)\theta = 2\cos\theta\cos n\theta$, they satisfy
\[
T_n(\cos\theta) = \cos(n\theta), \qquad T_n\left(\frac{z + z^{-1}}{2}\right) = \frac{z^n + z^{-n}}{2}, \qquad T_m(T_n(x)) = T_{mn}(x).
\]
\end{definition}

\begin{problem}
Factor $a^3 + b^3 + c^3 - 3abc$, deduce that $a+b+c \mid a^3+b^3+c^3 - 3abc$ for all integers $a,b,c$, and characterize when $a^3+b^3+c^3 = 3abc$ for real $a,b,c$.
\end{problem}
\begin{solution}
Viewing $P(a) = a^3+b^3+c^3-3abc$ as a polynomial in $a$, substitute
$a = -(b+c)$: we get $-(b+c)^3 + b^3 + c^3 + 3bc(b+c) = 0$. Hence
$(a+b+c) \mid P$, and polynomial division yields
\[
a^3+b^3+c^3-3abc = (a+b+c)(a^2+b^2+c^2-ab-bc-ca) = \frac{1}{2}(a+b+c)\bigl((a-b)^2+(b-c)^2+(c-a)^2\bigr).
\]
This gives the divisibility over $\ZZ$, and shows $a^3+b^3+c^3=3abc \iff a+b+c=0$ or $a=b=c$ over $\RR$.
\end{solution}

\section{Complex Numbers and Trigonometry}

\begin{definition}
$\CC = \{a+bi : a,b\in\RR\}$, $i^2=-1$. For $z = a+bi$: the conjugate is
$\bar z = a - bi$, the modulus $|z| = \sqrt{a^2+b^2} = \sqrt{z\bar z}$, and the
polar form is $z = r(\cos\theta + i\sin\theta) = re^{i\theta}$.
\end{definition}

\begin{theorem}[Euler / de Moivre and Triangle Identities]
\leavevmode
\begin{enumerate}[label=(\alph*)]
\item $e^{i\theta} = \cos\theta + i \sin\theta$, hence
$(\cos\theta + i\sin\theta)^n = \cos n\theta + i \sin n\theta$.
\item If $A + B + C = \pi$, then:
\begin{align*}
\tan A + \tan B + \tan C &= \tan A \tan B \tan C, &
\sin A + \sin B + \sin C &= 4\cos\frac{A}{2}\cos\frac{B}{2}\cos\frac{C}{2},\\
\cot\frac{A}{2} + \cot\frac{B}{2} + \cot\frac{C}{2} &= \cot\frac{A}{2}\cot\frac{B}{2}\cot\frac{C}{2}, &
\cos A + \cos B + \cos C &= 1 + 4\sin\frac{A}{2}\sin\frac{B}{2}\sin\frac{C}{2}.
\end{align*}
\end{enumerate}
\end{theorem}
\begin{proof}
De Moivre follows by induction from
$(\cos\alpha + i \sin\alpha)(\cos\beta + i\sin\beta) = \cos(\alpha+\beta) + i\sin(\alpha+\beta)$.
For $\tan$: $\tan(A+B) = -\tan C \implies \frac{\tan A+\tan B}{1-\tan A\tan B} = -\tan C$, which rearranges to the identity.
\end{proof}

\textbf{Roots of unity.} The solutions of $z^n=1$ are
$\omega^k = e^{2\pi i k/n}$, $k=0,\dots,n-1$. They satisfy
\[
z^n - 1 = \prod_{k=0}^{n-1}(z-\omega^k), \qquad
\sum_{k=0}^{n-1} \omega^{km} =
\begin{cases} n & n \mid m,\\ 0 & n \nmid m,\end{cases}
\]
the latter being the \textbf{roots of unity filter}, which extracts every $n$-th coefficient of $P(x) = \sum a_m x^m$:
\[
\sum_{m \equiv r \pmod n} a_m = \frac{1}{n}\sum_{k=0}^{n-1} \omega^{-kr} P(\omega^k).
\]

\begin{example}
Evaluate $\cos\frac{2\pi}{5} + \cos\frac{4\pi}{5}$.
\end{example}
\begin{solution}
The five 5th roots of unity sum to $0$. Taking real parts:
$1 + 2\cos\frac{2\pi}{5} + 2\cos\frac{4\pi}{5} = 0$, so the sum equals $-\tfrac12$.
\end{solution}

\begin{problem}
Prove that $\displaystyle\sum_{k \equiv 0 \pmod 3} \binom{n}{k}
= \frac{2^n + 2\cos\frac{n\pi}{3}}{3}$ and $\displaystyle\prod_{k=1}^{n-1}\sin\frac{k\pi}{n} = \frac{n}{2^{n-1}}$.
\end{problem}
\begin{solution}
Let $\omega = e^{2\pi i/3}$. By the roots of unity filter applied to
$(1+x)^n = \sum_k \binom nk x^k$,
\[
\sum_{3 \mid k}\binom nk = \frac{(1+1)^n + (1+\omega)^n + (1+\omega^2)^n}{3}.
\]
Now $1+\omega = e^{i\pi/3}$ and $1+\omega^2 = e^{-i\pi/3}$, so
$(1+\omega)^n + (1+\omega^2)^n = 2\cos\frac{n\pi}{3}$, giving the first claim.
For the product, divide $z^n - 1 = \prod_{k=0}^{n-1}(z - e^{2\pi i k/n})$ by $z-1$, set $z=1$, and take moduli using $|1 - e^{2i\theta}| = 2|\sin\theta|$.
\end{solution}

\section{Inequalities}

\subsection{AM--GM and Cauchy--Schwarz}
\begin{theorem}[AM--GM]
For nonnegative reals $a_1,\dots,a_n$,
\[
\frac{a_1+\dots+a_n}{n} \ge \sqrt[n]{a_1\cdots a_n},
\]
with equality iff all $a_i$ are equal.
\end{theorem}

\begin{proof}[Proof (Cauchy's forward--backward induction)]
\emph{Case $n=2$:} $(\sqrt{a_1}-\sqrt{a_2})^2 \ge 0$.
\emph{Powers of 2:} if the claim holds for $n$, apply it to two blocks of $n$
and then the case $n=2$ to their averages; this proves it for $2n$, hence all
$n = 2^k$.
\emph{Backward step:} if the claim holds for $n$, apply it for $n$ to
$a_1,\dots,a_{n-1}, A$ where $A = \frac{a_1+\dots+a_{n-1}}{n-1}$. The
arithmetic mean of these $n$ numbers is again $A$, so
$A \ge (a_1\cdots a_{n-1} A)^{1/n}$, i.e.\ $A^{n-1} \ge a_1\cdots a_{n-1}$,
which is the claim for $n-1$. Combining, AM--GM holds for all $n$. Equality
tracing shows all terms equal.
\end{proof}

\begin{theorem}[Cauchy--Schwarz]
For reals $a_i, b_i$,
\[
\left(\sum_{i=1}^n a_i b_i\right)^2 \le
\left(\sum_{i=1}^n a_i^2\right)\left(\sum_{i=1}^n b_i^2\right),
\]
with equality iff the vectors $(a_i)$ and $(b_i)$ are proportional.
\end{theorem}
\begin{proof}
The quadratic $q(t) = \sum_i (a_i t + b_i)^2 = t^2\sum a_i^2 + 2t \sum a_ib_i +
\sum b_i^2$ is $\ge 0$ for all $t$, so its discriminant is $\le 0$, which is the
inequality. Equality means $q$ has a real root $t_0$, i.e.\ $b_i = -t_0 a_i$ for all $i$.
\end{proof}

\begin{corollary}[Cauchy--Schwarz in Engel form / Titu's Lemma]
For reals $x_i$ and positive $y_i$:
$\displaystyle \sum \frac{x_i^2}{y_i} \ge \frac{(\sum x_i)^2}{\sum y_i}$.
\end{corollary}
\begin{proof}
Apply Cauchy--Schwarz to $a_i = x_i/\sqrt{y_i}$, $b_i = \sqrt{y_i}$.
\end{proof}

\begin{example}[Nesbitt's Inequality]
For $a,b,c>0$:
$\dfrac{a}{b+c}+\dfrac{b}{c+a}+\dfrac{c}{a+b} \ge \dfrac32$.
\end{example}
\begin{solution}
By Engel form,
\[
\sum \frac{a}{b+c} = \sum \frac{a^2}{a(b+c)}
\ge \frac{(a+b+c)^2}{2(ab+bc+ca)} \ge \frac{3}{2},
\]
the last step because $(a+b+c)^2 \ge 3(ab+bc+ca)$, i.e.\
$a^2+b^2+c^2 \ge ab+bc+ca$, which follows from summing
$\frac12\bigl((a-b)^2 + (b-c)^2 + (c-a)^2\bigr) \ge 0$.
\end{solution}

\subsection{Rearrangement, Chebyshev, and Schur's Inequalities}

\begin{theorem}[Rearrangement and Chebyshev Inequalities]
Let $a_1 \le a_2 \le \dots \le a_n$ and $b_1 \le b_2 \le \dots \le b_n$ be real numbers. For any permutation $\sigma$ of $\{1,\dots,n\}$,
\[
\sum_{i=1}^n a_i b_{n+1-i} \le \sum_{i=1}^n a_i b_{\sigma(i)} \le \sum_{i=1}^n a_i b_i,
\qquad\text{and}\qquad
\frac{1}{n}\sum_{i=1}^n a_i b_i \ge \left(\frac{1}{n}\sum_{i=1}^n a_i\right)\left(\frac{1}{n}\sum_{i=1}^n b_i\right).
\]
\end{theorem}
\begin{proof}
If $\sigma$ has an inversion $i < j$ with $\sigma(i) > \sigma(j)$, swapping $\sigma(i)$ and $\sigma(j)$ changes the sum by $(a_j - a_i)(b_{\sigma(i)} - b_{\sigma(j)}) \ge 0$. Uncrossing all inversions reaches $\sigma = \mathrm{id}$. Averaging the $n$ cyclic shifts gives Chebyshev's inequality.
\end{proof}

\begin{theorem}[Schur's Inequality]
For nonnegative reals $x, y, z$ and $r > 0$,
\[
x^r(x-y)(x-z) + y^r(y-z)(y-x) + z^r(z-x)(z-y) \ge 0,
\]
with equality iff $x=y=z$ or two variables are equal and the third is $0$.
For $r = 1$: $x^3+y^3+z^3+3xyz \ge \sum_{\text{sym}} x^2 y$, or $p^3 - 4pq + 9r \ge 0$ ($p=\sum x, q=\sum xy, r=xyz$).
\end{theorem}
\begin{proof}
WLOG assume $x \ge y \ge z \ge 0$. Rewrite the left-hand side as
\[
(x-y)\bigl(x^r(x-z) - y^r(y-z)\bigr) + z^r(z-x)(z-y).
\]
Since $x^r \ge y^r \ge 0$ and $x-z \ge y-z \ge 0$, the first term is $\ge 0$; since $z-x \le 0$ and $z-y \le 0$, the second term is also $\ge 0$.
\end{proof}

\begin{remark}[Homogenization and Ravi Substitution]
Given a constraint like $a+b+c=1$, multiply terms by powers of $a+b+c$ to make the inequality homogeneous. When $a,b,c$ are the sides of a triangle, the \textbf{Ravi substitution} $a = y+z, b = z+x, c = x+y$ ($x,y,z > 0$) eliminates the triangle inequality constraints.
\end{remark}

\begin{problem}[IMO 2000/2]
Let $a,b,c > 0$ with $abc = 1$. Prove that
$$\left(a - 1 + \frac{1}{b}\right)\left(b - 1 + \frac{1}{c}\right)\left(c - 1 + \frac{1}{a}\right) \le 1$$
\end{problem}
\begin{solution}
Substitute $a = x/y, b = y/z, c = z/x$ for $x,y,z > 0$. Clearing denominators, the inequality becomes
$(x+y-z)(y+z-x)(z+x-y) \le xyz$. At most one factor on the left is $\le 0$ (as the sum of any two is positive); if one is $\le 0$, the product is $\le 0 < xyz$. Otherwise expanding gives $x^3+y^3+z^3+3xyz \ge \sum_{\text{sym}}x^2y$, which is Schur for $r=1$.
\end{solution}

\subsection{Jensen's, Hölder's, and Minkowski's Inequalities [Advanced]}
\begin{definition}
$f$ is \textbf{convex} on an interval $I$ if
$f(\lambda x + (1-\lambda)y) \le \lambda f(x) + (1-\lambda) f(y)$ for all
$x,y \in I$, $\lambda \in [0,1]$; \textbf{concave} if the inequality reverses.
If $f'' \ge 0$ on $I$ then $f$ is convex.
\end{definition}

\begin{theorem}[Jensen]
If $f$ is convex on $I$, $x_i \in I$, and $w_i > 0$ with $\sum w_i = 1$, then
\[
f\Bigl(\sum_{i=1}^n w_i x_i\Bigr) \le \sum_{i=1}^n w_i f(x_i).
\]
\end{theorem}
\begin{proof}
Induct on $n$; the case $n=2$ is the definition. For $n>2$ write
$\sum_{i=1}^{n} w_i x_i = w_n x_n + (1-w_n) y$ where
$y = \sum_{i<n} \frac{w_i}{1-w_n} x_i \in I$. Then by convexity and the
induction hypothesis,
\[
f\Bigl(\sum w_i x_i\Bigr) \le w_n f(x_n) + (1-w_n) f(y)
\le w_n f(x_n) + (1-w_n)\sum_{i<n}\tfrac{w_i}{1-w_n}f(x_i)
= \sum w_i f(x_i). \qedhere
\]
\end{proof}

\begin{theorem}[Hölder and Minkowski]
For positive reals and exponents $p, q > 1$ with $\frac1p + \frac1q = 1$:
\[
\sum_{i=1}^n a_i b_i \le \Bigl(\sum a_i^p\Bigr)^{1/p}\Bigl(\sum b_i^q\Bigr)^{1/q}, \qquad
\Bigl(\sum (a_i+b_i)^p\Bigr)^{1/p} \le \Bigl(\sum a_i^p\Bigr)^{1/p} + \Bigl(\sum b_i^p\Bigr)^{1/p}.
\]
In three-sequence olympiad form: $\bigl(\sum a_i\bigr)\bigl(\sum b_i\bigr)\bigl(\sum c_i\bigr) \ge \bigl(\sum \sqrt[3]{a_i b_i c_i}\bigr)^3$.
\end{theorem}
\begin{proof}
For Hölder, normalize $\sum a_i^p = \sum b_i^q = 1$. By weighted AM--GM,
$a_i b_i = (a_i^p)^{1/p}(b_i^q)^{1/q} \le \frac{a_i^p}{p} + \frac{b_i^q}{q}$;
summing gives $\sum a_ib_i \le \frac1p + \frac1q = 1$. For Minkowski, let $s_i = a_i+b_i$; by Hölder,
$\sum s_i^p = \sum a_i s_i^{p-1} + \sum b_i s_i^{p-1} \le \bigl[(\sum a_i^p)^{1/p} + (\sum b_i^p)^{1/p}\bigr](\sum s_i^p)^{1/q}$, and divide by $(\sum s_i^p)^{1/q}$.
\end{proof}

\subsection{Muirhead, Karamata, SOS, and the \texorpdfstring{$uvw$}{uvw} Method [Advanced]}

\begin{theorem}[Muirhead, Karamata, SOS, and Teichmüller's $uvw$ Theorem]
\leavevmode
\begin{enumerate}[label=(\alph*)]
\item \textbf{(Muirhead \& Karamata)} Say $\alpha \succ \beta$ if $\alpha_1\ge\dots\ge\alpha_n$, $\beta_1\ge\dots\ge\beta_n$, $\sum_{i=1}^k \alpha_i \ge \sum_{i=1}^k \beta_i$ ($k<n$), and $\sum_{i=1}^n \alpha_i = \sum_{i=1}^n \beta_i$. If $\alpha \succ \beta$, then $\sum_{\sigma\in S_n} x_{\sigma(1)}^{\alpha_1}\cdots x_{\sigma(n)}^{\alpha_n} \ge \sum_{\sigma\in S_n} x_{\sigma(1)}^{\beta_1}\cdots x_{\sigma(n)}^{\beta_n}$ for $x_i > 0$, and $\sum_{i=1}^n f(\alpha_i) \ge \sum_{i=1}^n f(\beta_i)$ for any convex $f$ (proved via Abel summation on the nonincreasing secant slopes $c_i = \frac{f(\alpha_i)-f(\beta_i)}{\alpha_i-\beta_i}$).
\item \textbf{(SOS Method)} $\sum_{\text{cyc}} S_a(b-c)^2 \ge 0$ holds if $S_a, S_b, S_c \ge 0$, or if $S_a, S_c \ge 0$ and $S_a + 2S_b \ge 0, S_c + 2S_b \ge 0$.
\item \textbf{($uvw$ / $pqr$ Theorem)} For $a,b,c \ge 0$ with fixed $p = a+b+c$ and $q = ab+bc+ca$, the product $r = abc$ attains its extrema only when \emph{two variables are equal} or \emph{one variable is $0$} (since $a,b,c$ are the real roots of $t^3 - pt^2 + qt = r$, and shifting the horizontal line $y=r$ reaches an extremum when the line is tangent to the cubic—a double root—or hits $r=0$). Thus any inequality $f(p,q) + Cr \ge 0$ linear in $r$ need only be checked for $a=b$ and $c=0$.
\end{enumerate}
\end{theorem}

\begin{problem}
Let $a,b,c \ge 0$ with $a+b+c = 3$.
(a) If $a,b,c > 0$, prove $\dfrac{1}{a^2} + \dfrac{1}{b^2} + \dfrac{1}{c^2} \ge a^2+b^2+c^2$.
(b) Prove $a^2+b^2+c^2+abc \ge 4$.
\end{problem}
\begin{solution}
(a) Use the \emph{tangent line trick} on $f(x) = \frac1{x^2}-x^2$ at $x = 1$:
$f(1)=0$, $f'(1) = -4$, and
\[
\frac{1}{x^2} - x^2 - \bigl(-4(x-1)\bigr) = \frac{(x-1)^2(-x^2+2x+1)}{x^2} \ge 0 \quad \text{for } x \in (0, 1+\sqrt{2}].
\]
If $a,b,c \le 1+\sqrt2$, summing gives $\sum f(a) \ge -4(a+b+c-3) = 0$. If say $c > 1+\sqrt2$, then $a+b < 2-\sqrt2$, so $\min(a,b) < 1-\frac{\sqrt2}{2} < 0.3$ and $\frac1{a^2}+\frac1{b^2} > \frac{1}{0.09} > 11 > 9 \ge a^2+b^2+c^2$.

(b) In $p,q,r$ form with $p=3$, $a^2+b^2+c^2+abc-4 = 5 - 2q + r \ge 0$, which is linear in $r$. By $uvw$: if $c=0$, $a^2+b^2 \ge \frac{(a+b)^2}{2} = \frac{9}{2} > 4$; if $a=b=x$ and $c=3-2x$ ($x\in[0,3/2]$), then $2x^2+(3-2x)^2+x^2(3-2x)-4 = (x-1)^2(5-2x) \ge 0$.
\end{solution}

\section{Floor Functions, Beatty's Theorem, and Sequence Asymptotics}

\begin{theorem}[Hermite's Identity, Beatty's Theorem, and Telescoping Asymptotics]
\leavevmode
\begin{enumerate}[label=(\alph*)]
\item \textbf{(Hermite)} For $x \in \RR$ and $n \in \NN$, $\sum_{k=0}^{n-1}\lfloor x + \frac{k}{n}\rfloor = \lfloor nx \rfloor$.
\item \textbf{(Beatty)} If $\alpha, \beta > 1$ are irrational with $\frac{1}{\alpha}+\frac{1}{\beta}=1$, then $A = \{\lfloor n\alpha\rfloor : n\in\NN\}$ and $B = \{\lfloor m\beta\rfloor : m\in\NN\}$ partition $\NN$.
\item \textbf{(Telescoping Powers)} For $x_{n+1} = x_n + c x_n^{-\alpha}$ ($c>0$), raising to power $\alpha+1$ makes $x_{n+1}^{\alpha+1} - x_n^{\alpha+1} \to (\alpha+1)c$, yielding $x_n \sim \bigl((\alpha+1)cn\bigr)^{1/(\alpha+1)}$.
\end{enumerate}
\end{theorem}
\begin{proof}[Proof of (a) and (b)]
(a) Write $x = m + r/n + \delta$ with $m=\lfloor x\rfloor$, $r\in\{0,\dots,n-1\}$, $0\le\delta < 1/n$; both sides equal $nm+r$.
(b) For any integer $k \ge 1$, $\{k/\alpha\} + \{k/\beta\} = 1$ since $k/\alpha + k/\beta = k \in \ZZ$ and both are irrational. Thus $\lfloor k/\alpha \rfloor + \lfloor k/\beta \rfloor = k-1$. Subtracting this for $k$ and $k+1$ gives $\bigl(\lfloor\frac{k+1}{\alpha}\rfloor - \lfloor\frac{k}{\alpha}\rfloor\bigr) + \bigl(\lfloor\frac{k+1}{\beta}\rfloor - \lfloor\frac{k}{\beta}\rfloor\bigr) = 1$, so $k$ lies in exactly one of $A, B$.
\end{proof}

\begin{problem}
(a) Let $x_1 = 1$ and $x_{n+1} = x_n + \frac{1}{x_n}$. Find $\lim_{n\to\infty}\frac{x_n}{\sqrt{n}}$.
(b) Prove that $\lfloor (2+\sqrt{3})^n \rfloor$ is odd for every positive integer $n$.
\end{problem}
\begin{solution}
(a) Squaring gives $x_{n+1}^2 - x_n^2 = 2 + x_n^{-2} > 2$, so summing from $1$ to $n-1$ gives $x_n^2 \ge 2n-1$. Substituting this back gives $x_n^2 = 2n-1 + \sum_{k=1}^{n-1}x_k^{-2} = 2n + O(\ln n)$, so $x_n/\sqrt{n} \to \sqrt{2}$.
(b) Let $u_n = (2+\sqrt{3})^n + (2-\sqrt{3})^n$. Then $u_0=2, u_1=4$, and $u_{n+1} = 4u_n - u_{n-1}$, so $u_n$ is always an even integer. Since $0 < (2-\sqrt{3})^n < 1$, we have $\lfloor (2+\sqrt{3})^n \rfloor = u_n - 1$, which is odd.
\end{solution}

\section{Generating Functions [Advanced]}

\begin{definition}
The \textbf{ordinary generating function} of a sequence $(a_n)_{n\ge0}$ is the
formal power series $A(x) = \sum_{n\ge0} a_n x^n$.
\end{definition}

Key dictionary: addition of series $\leftrightarrow$ termwise addition;
multiplication $\leftrightarrow$ convolution \\
$\bigl(\sum a_nx^n\bigr)\bigl(\sum b_nx^n\bigr) = \sum_n \bigl(\sum_{k} a_k b_{n-k}\bigr)x^n$;
$\frac{1}{1-x} = \sum x^n$; and the generalized binomial theorem
$\frac{1}{(1-x)^k} = \sum_{n\ge0}\binom{n+k-1}{k-1}x^n$.

\begin{example}[Fibonacci Closed Form]
Let $F_0=0$, $F_1=1$, $F_{n+2}=F_{n+1}+F_n$. Then with
$\varphi = \frac{1+\sqrt5}{2}$, $\psi = \frac{1-\sqrt5}{2}$:
$F_n = \frac{\varphi^n - \psi^n}{\sqrt5}$.
\end{example}
\begin{solution}
Let $F(x)=\sum F_n x^n$. The recurrence gives
$F(x) - x = x F(x) + x^2 F(x)$, so $F(x) = \frac{x}{1-x-x^2} = \frac{1}{\sqrt5}\left(\frac{1}{1-\varphi x} - \frac{1}{1-\psi x}\right)$,
and expanding each geometric series yields the formula.
\end{solution}

\begin{problem}[Sicherman Dice]
Show that one cannot relabel two six-sided dice with positive integers so that
every pairwise sum $2,\dots,12$ is equally likely, but find the unique nonstandard positive-integer labeling (Sicherman dice) giving the same sum distribution as two standard dice.
\end{problem}
\begin{solution}
A die with faces $a_1,\dots,a_6$ has generating function $D(x)=\sum x^{a_i}$. A uniform distribution on the $11$ sums $\{2,\dots,12\}$ is impossible since $11 \nmid 36 = D_1(1)D_2(1)$. For standard dice:
\[
\bigl(x + \dots + x^6\bigr)^2 = x^2(1+x)^2(1+x+x^2)^2(1-x+x^2)^2.
\]
We must split these irreducible factors into $D_1(x)D_2(x)$ with $D_i(1)=6$ and $D_i(0)=0$. Since $(1+x)(1+x+x^2)\big|_{x=1}=6$ and $(1-x+x^2)\big|_{x=1}=1$, each die must receive $x(1+x)(1+x+x^2)$. Giving both factors of $(1-x+x^2)$ to $D_2$ yields
$$D_1(x) = x(1+x)(1+x+x^2) = x+2x^2+2x^3+x^4$$
$$D_2(x) = x(1+x)(1+x+x^2)(1-x+x^2)^2 = x+x^3+x^4+x^5+x^6+x^8$$
corresponding to faces $\{1,2,2,3,3,4\}$ and $\{1,3,4,5,6,8\}$.
\end{solution}

\section{Calculus in Olympiads [Advanced]}

\begin{theorem}[Key Facts]
\leavevmode
\begin{enumerate}[label=(\alph*)]
\item If $f'(x) \ge 0$ on an interval, $f$ is nondecreasing there.
\item \textbf{(Rolle)} If $f$ is differentiable and $f(a)=f(b)$, some $c\in(a,b)$ has
$f'(c)=0$. Consequently, between two real roots of a polynomial lies a root of
its derivative; a polynomial with all real roots has all derivatives with all
real roots.
\item If $f''\ge 0$, then $f$ is convex and lies above each tangent line:
$f(x) \ge f(a) + f'(a)(x-a)$ (the \emph{tangent line trick}).
\item \textbf{(Smoothing)} To extremize a symmetric expression under a constraint,
examine what happens when two variables move toward (or away from) each other
with their sum or product fixed.
\end{enumerate}
\end{theorem}
\begin{proof}[Proof of (c)]
Let $g(x) = f(x)-f(a)-f'(a)(x-a)$. Then $g(a)=0$, $g'(a)=0$, and $g''=f''\ge0$
means $g'$ is nondecreasing; so $g' \le 0$ before $a$ and $g'\ge 0$ after,
hence $g$ attains its minimum $0$ at $a$.
\end{proof}

\begin{problem}
(a) Let $x,y,z>0$ with $x+y+z=1$. Find the minimum of $\sum \frac{x}{1-x}$.
(b) If $P(x) = x^n + a_{n-1}x^{n-1} + \dots + a_0$ ($n\ge 2$) has $n$ real roots, prove $(n-1)a_{n-1}^2 \ge 2na_{n-2}$.
\end{problem}
\begin{solution}
(a) $f(t) = \frac{t}{1-t}$ on $(0,1)$ has $f''(t) = \frac{2}{(1-t)^3} > 0$. By Jensen, $\sum f(x_i) \ge 3f(1/3) = \frac{3}{2}$.
(b) By Rolle's theorem applied $n-2$ times, $\frac{P^{(n-2)}(x)}{(n-2)!} = \frac{n(n-1)}{2}x^2 + (n-1)a_{n-1}x + a_{n-2}$ has two real roots, so its discriminant $(n-1)^2a_{n-1}^2 - 2n(n-1)a_{n-2} \ge 0$.
\end{solution}

%======================================================================
\part{Combinatorics}
%======================================================================

\section{Basic Counting and Binomial Identities}

\textbf{Fundamental principles.} If a procedure has $m$ outcomes and an
independent second stage has $n$, together there are $mn$ (product rule);
disjoint cases add (sum rule).

\begin{definition}
$n! = n(n-1)\cdots 1$ counts orderings of $n$ objects;
$\binom{n}{k} = \frac{n!}{k!(n-k)!}$ counts $k$-element subsets of an
$n$-element set; $2^n$ counts all subsets.
\end{definition}

\begin{theorem}[Binomial Theorem]
$(x+y)^n = \sum_{k=0}^n \binom nk x^k y^{n-k}$.
\end{theorem}
\begin{proof}
Expanding the product of $n$ brackets, each term chooses $x$ from $k$ brackets
and $y$ from the rest; there are $\binom nk$ ways to choose which.
\end{proof}

\textbf{Core identities} (each with a combinatorial proof worth internalizing):
\begin{align*}
\binom{n}{k} &= \binom{n}{n-k}, &
\binom{n}{k} &= \binom{n-1}{k}+\binom{n-1}{k-1} \quad (\text{Pascal}),\\
k\binom{n}{k} &= n\binom{n-1}{k-1} \quad (\text{Leader}), &
\sum_{k=r}^n \binom{k}{r} &= \binom{n+1}{r+1} \quad (\text{Hockey-Stick}),\\
\sum_{k=0}^r\binom{m}{k}\binom{n}{r-k} &= \binom{m+n}{r} \quad (\text{Vandermonde}), &
\#\{x_1+\dots+x_k &= n,\ x_i \ge 0\} = \binom{n+k-1}{k-1} \quad (\text{Stars \& Bars}).
\end{align*}

\begin{example}
How many lattice paths go from $(0,0)$ to $(m,n)$ using unit steps right and up?
\end{example}
\begin{solution}
A path is a sequence of $m+n$ steps of which $m$ are ``right'':
$\binom{m+n}{m}$ paths.
\end{solution}

\begin{problem}
Prove $\displaystyle\sum_{k=0}^{n} \binom{n}{k}^2 = \binom{2n}{n}$.
\end{problem}
\begin{solution}
Count committees of $n$ people chosen from $n$ men and $n$ women. Choosing
$k$ men and $n-k$ women can be done in $\binom nk\binom{n}{n-k} = \binom nk^2$
ways; summing over $k$ counts all $\binom{2n}{n}$ committees.
\end{solution}

\section{Mathematical Induction}

\begin{theorem}[Principle of Induction]
Let $S \subseteq \NN$ contain $1$, and suppose $n \in S \implies n+1 \in S$.
Then $S = \NN$. Equivalent forms: \emph{strong induction} (assume all of
$1,\dots,n$ to prove $n+1$) and \emph{well-ordering} (every nonempty subset of
$\NN$ has a least element).
\end{theorem}

\begin{example}
Prove that $\sum_{k=1}^{n} k^3 = \left(\frac{n(n+1)}{2}\right)^2$.
\end{example}
\begin{solution}
Base $n=1$: both sides equal $1$. Step: assuming the identity for $n$,
\[
\sum_{k=1}^{n+1}k^3 = \frac{n^2(n+1)^2}{4} + (n+1)^3
= \frac{(n+1)^2\bigl(n^2+4n+4\bigr)}{4} = \left(\frac{(n+1)(n+2)}{2}\right)^2.
\qedhere
\]
\end{solution}

\begin{problem}[Strong Induction]
Every integer $n \ge 2$ is a product of primes.
\end{problem}
\begin{solution}
Strong induction. If $n$ is prime, done. Otherwise $n = ab$ with
$2\le a,b < n$; by the strong induction hypothesis each of $a, b$ is a product
of primes, hence so is $n$.
\end{solution}

\begin{remark}
Powerful variants: strengthen the inductive hypothesis, induct on an unusual quantity (sum of coordinates, number
of inversions), \emph{downward} induction, and the extremal principle
(consider a minimal counterexample and derive a contradiction).
\end{remark}

\section{Recurrence Relations, Catalan Numbers, and Partitions}

\begin{theorem}[Linear Recurrences with Constant Coefficients]
Suppose $a_{n} = c_1 a_{n-1} + \dots + c_k a_{n-k}$ and the characteristic
polynomial $x^k - c_1x^{k-1} - \dots - c_k$ has distinct roots
$r_1,\dots,r_k$. Then $a_n = \sum_i \alpha_i r_i^n$ for constants $\alpha_i$
determined by initial values. (A root $r$ of multiplicity $m$ contributes
$(\beta_0 + \beta_1 n + \dots + \beta_{m-1}n^{m-1})r^n$.)
\end{theorem}
\begin{proof}[Proof sketch]
Each $r_i^n$ satisfies the recurrence by definition of the characteristic
polynomial; solutions form a $k$-dimensional vector space (a solution is
determined by $k$ initial values); the $k$ sequences $(r_i^n)$ are linearly
independent (Vandermonde determinant), hence a basis.
\end{proof}

\begin{example}
Count binary strings of length $n$ with no two consecutive $1$s.
\end{example}
\begin{solution}
Let $a_n$ be the count. A valid string ends in $0$ (preceded by any valid
string of length $n-1$) or in $01$ (preceded by any valid string of length
$n-2$): $a_n = a_{n-1}+a_{n-2}$, with $a_1 = 2$, $a_2 = 3$. Thus
$a_n = F_{n+2}$ (Fibonacci).
\end{solution}

\subsection{Catalan Numbers, Partitions, and Non-Intersecting Paths}

\begin{theorem}[Catalan Numbers]
The $n$-th \textbf{Catalan number} $C_n = \frac{1}{n+1}\binom{2n}{n}$ counts:
(1) \textbf{Dyck paths} from $(0,0)$ to $(n,n)$ with steps $(1,0)$ and $(0,1)$ never rising above $y=x$;
(2) balanced strings of $n$ pairs of parentheses;
(3) triangulations of a convex $(n+2)$-gon;
and satisfies $C_{n+1} = \sum_{k=0}^n C_k C_{n-k}$ ($C_0=1$).
\end{theorem}
\begin{proof}[Proof via André's Reflection Principle]
There are $\binom{2n}{n}$ total paths from $(0,0)$ to $(n,n)$. A path is \emph{bad} if it crosses above $y=x$, so it touches $y = x+1$. Reflect the portion of a bad path \emph{after} its first touch of $y = x+1$ across the line $y = x+1$. Since $(n,n)$ reflects to $(n-1,n+1)$, this gives a bijection between bad paths and all paths from $(0,0)$ to $(n-1,n+1)$, of which there are $\binom{2n}{n-1}$. Thus
\[
C_n = \binom{2n}{n} - \binom{2n}{n-1} = \frac{1}{n+1}\binom{2n}{n}. \qedhere
\]
\end{proof}

\begin{theorem}[Euler's Partition Identity and LGV Lemma]
\leavevmode
\begin{enumerate}[label=(\alph*)]
\item \textbf{(Euler)} For every $n \ge 1$, the number of partitions of $n$ into \emph{odd} parts equals the number of partitions of $n$ into \emph{distinct} parts: $\prod_{k\ge 1}\frac{1}{1-x^{2k-1}} = \prod_{k\ge 1}\frac{1-x^{2k}}{1-x^k} = \prod_{k\ge 1}(1+x^k)$.
\item \textbf{(Lindström--Gessel--Viennot)} In a planar DAG where any path $A_i \to B_j$ and $A_{i'} \to B_{j'}$ ($i<i', j>j'$) must cross, the number of vertex-disjoint path systems from $(A_1,\dots,A_k)$ to $(B_1,\dots,B_k)$ equals $\det\bigl(e(A_i, B_j)\bigr)_{i,j=1}^k$.
\end{enumerate}
\end{theorem}

\begin{problem}
Let $a_1 = 1$ and $a_{n+1} = 2a_n + 1$. Find a closed form and prove it.
\end{problem}
\begin{solution}
Add 1 to both sides: $a_{n+1}+1 = 2(a_n + 1)$. So $b_n = a_n+1$ is geometric
with $b_1 = 2$: $b_n = 2^n$, hence $a_n = 2^n - 1$. (Induction verifies:
$a_{n+1} = 2(2^n-1)+1 = 2^{n+1}-1$.)
\end{solution}

\section{The Pigeonhole Principle}

\begin{theorem}[Pigeonhole]
If $n+1$ objects are placed into $n$ boxes, some box contains at least two
objects. More generally, if $N$ objects go into $n$ boxes, some box contains
at least $\lceil N/n \rceil$ objects.
\end{theorem}
\begin{proof}
If every box had at most $\lceil N/n\rceil - 1 < N/n$ objects, the total would
be less than $N$.
\end{proof}

\begin{example}
Among any $n+1$ integers chosen from $\{1, 2, \dots, 2n\}$, one divides another.
\end{example}
\begin{solution}
Write each chosen integer as $2^a \cdot m$ with $m$ odd. There are only $n$
odd numbers in $\{1,\dots,2n\}$, so two chosen integers share the same odd
part $m$: say $2^a m$ and $2^b m$ with $a<b$. Then the first divides the second.
\end{solution}

\begin{problem}
Prove that among any $n$ integers, some nonempty subset has sum divisible by $n$.
\end{problem}
\begin{solution}
Let the integers be $a_1,\dots,a_n$ and consider the prefix sums
$s_0 = 0,\ s_1 = a_1,\ \dots,\ s_n = a_1+\dots+a_n$. These are $n+1$ numbers
and only $n$ residues mod $n$, so $s_i \equiv s_j \pmod n$ for some $i<j$;
then $a_{i+1}+\dots+a_j = s_j - s_i$ is divisible by $n$.
\end{solution}

\section{Bijective Counting, Inclusion--Exclusion, and Posets}

\begin{definition}
For finite sets, $|A \cup B| = |A|+|B|-|A\cap B|$; more generally,
\textbf{inclusion--exclusion}:
\[
\Bigl|\bigcup_{i=1}^n A_i\Bigr| = \sum_{\emptyset \ne S \subseteq [n]}
(-1)^{|S|+1}\Bigl|\bigcap_{i\in S}A_i\Bigr|.
\]
\end{definition}
\begin{proof}
Fix an element $x$ lying in exactly $k \ge 1$ of the sets. Its contribution to
the right-hand side is $\sum_{j=1}^{k}(-1)^{j+1}\binom kj = 1 - (1-1)^k = 1$,
matching the left side.
\end{proof}

\textbf{Bijections and double counting.} If a bijection exists between two
finite sets, they have equal size; counting one quantity in two ways yields an
identity. Injections give $|A|\le|B|$.

\begin{example}[Derangements]
The number of permutations of $[n]$ with no fixed point is
$D_n = n!\sum_{k=0}^{n}\frac{(-1)^k}{k!}$.
\end{example}
\begin{solution}
Let $A_i$ be the set of permutations fixing $i$. By inclusion--exclusion,
$$|\bigcup A_i| = \sum_k (-1)^{k+1}\binom nk (n-k)!$$so
$D_n = n! - \sum_{k\ge1}(-1)^{k+1}\frac{n!}{k!} = n!\sum_{k=0}^n \frac{(-1)^k}{k!}$.
\end{solution}

\subsection{Sperner's Theorem and Dilworth's Theorem [Advanced]}

\begin{theorem}[LYM Inequality, Sperner, Mirsky, and Dilworth]
\leavevmode
\begin{enumerate}[label=(\alph*)]
\item \textbf{(LYM \& Sperner)} Let $\mathcal{F}$ be an antichain of subsets of $[n]$ (no set contains another). Then
\[
\sum_{A \in \mathcal{F}} \frac{1}{\binom{n}{|A|}} \le 1, \qquad\text{and hence}\qquad
|\mathcal{F}| \le \binom{n}{\lfloor n/2 \rfloor}.
\]
\item \textbf{(Mirsky \& Dilworth)} In any finite poset $(P, \le)$, the maximum size of a chain equals the minimum number of antichains needed to partition $P$, and the maximum size of an antichain equals the minimum number of chains needed to partition $P$.
\end{enumerate}
\end{theorem}
\begin{proof}
(a) Among all $n!$ maximal chains $\emptyset = C_0 \subset C_1 \subset \dots \subset C_n = [n]$, each $A \in \mathcal{F}$ belongs to $|A|!(n-|A|)!$ chains, and no chain contains two members of $\mathcal{F}$. Thus $\sum_{A \in \mathcal{F}} |A|!(n-|A|)! \le n!$, which divides by $n!$ to give LYM; using $\binom{n}{|A|} \le \binom{n}{\lfloor n/2\rfloor}$ gives Sperner's bound.
(b) For Mirsky, let $h(x)$ be the length of the longest chain ending at $x$; then $x < y \implies h(x) < h(y)$, so the level sets $A_k = \{x : h(x)=k\}$ partition $P$ into $\max h$ antichains.
\end{proof}

\begin{problem}[Double Counting and Erdős--Szekeres]
(a) At a party each pair of people are friends or strangers. Prove that the number
of people with an odd number of friends is even.
(b) Prove that any sequence of $mn+1$ real numbers has a nondecreasing subsequence of length $m+1$ or a decreasing subsequence of length $n+1$.
\end{problem}
\begin{solution}
(a) Count pairs (person, friend of that person): the total is
$\sum_v \deg(v) = 2E$, twice the number of friendships, hence even. A sum of
integers is even only if the number of odd summands is even.
(b) Map each $i \in \{1,\dots,mn+1\}$ to $(\ell_i, d_i)$, the lengths of the longest nondecreasing subsequence ending at $a_i$ and longest decreasing starting at $a_i$. For $i < j$, $a_i \le a_j \implies \ell_j \ge \ell_i+1$ and $a_i > a_j \implies d_i \ge d_j+1$, so the map is injective and cannot fit in $\{1,\dots,m\}\times\{1,\dots,n\}$.
\end{solution}

\section{Invariants, Monovariants, and Coloring Arguments}

\begin{definition}
Under a set of allowed moves on a state space:
\begin{itemize}
\item An \textbf{invariant} $I(S)$ is unchanged by every move ($I(S')=I(S)$). If $I(S_{\text{start}}) \ne I(S_{\text{target}})$, then $S_{\text{target}}$ is unreachable.
\item A \textbf{monovariant} $M(S)$ strictly increases (or strictly decreases) on every move. If $M$ takes values in a finite or discrete bounded set, the process must terminate.
\end{itemize}
\end{definition}

\begin{example}[Checkerboard Coloring]
An $8 \times 8$ board with two opposite corners removed cannot be tiled by $31$ dominoes of size $1 \times 2$.
\end{example}
\begin{solution}
Each $1 \times 2$ domino covers one black and one white square on a standard chessboard coloring, whereas the mutilated board has $32$ squares of one color and $30$ of the other.
\end{solution}

\begin{problem}[Grid Infection Monovariant]
Initially, $k$ unit squares of an $n \times n$ grid are infected. At each step, any uninfected square sharing an edge with at least $2$ infected squares becomes infected. Prove that if the entire $n \times n$ grid becomes infected, $k \ge n$.
\end{problem}
\begin{solution}
Let $P$ be the perimeter of the infected region. When a new square is infected, it shares $c \ge 2$ edges with already-infected squares, so those $c$ edges disappear from the boundary while $4-c$ edges are added. Thus $\Delta P = (4-c) - c = 4 - 2c \le 0$: the perimeter $P$ never increases! Initially $P_{\text{init}} \le 4k$, and at the end $P_{\text{final}} = 4n$. Hence $4n \le 4k \implies k \ge n$.
\end{solution}

\section{The Extremal Principle and Combinatorial Games}

\begin{theorem}[Sylvester--Gallai Theorem]
Given a finite set $S$ of at least $3$ non-collinear points in the plane, there exists a line passing through \emph{exactly two} points of $S$.
\end{theorem}
\begin{proof}[Proof (Kelly)]
Among all pairs $(P, \ell)$ where $\ell$ is a line through at least two points of $S$ and $P \in S \setminus \ell$, choose one minimizing the perpendicular distance $d(P, \ell) > 0$. If $\ell$ contained $\ge 3$ points of $S$, at least two of them ($B$ and $C$) would lie on the same side of the foot $H$ of $P$ on $\ell$, with $B$ between $H$ and $C$ (possibly $B=H$). Then the distance from $B$ to line $PC$ is strictly less than $PH = d(P, \ell)$, contradicting minimality.
\end{proof}

\begin{theorem}[Combinatorial Games, Nim, and Sprague--Grundy]
In a finite, impartial, normal-play game (last player to move wins):
\begin{itemize}
\item A position is a \textbf{losing ($P$) position} if all moves from it lead to $N$-positions, and a \textbf{winning ($N$) position} if some move leads to a $P$-position.
\item \textbf{(Bouton / Nim)} In Nim with pile sizes $(n_1,\dots,n_k)$, a state is a $P$-position iff the XOR sum $n_1 \oplus \dots \oplus n_k = 0$.
\item \textbf{(Sprague--Grundy)} Every impartial game position $G$ is equivalent to a single Nim pile of size $g(G) = \operatorname{mex}\{g(G') : G \to G'\}$, and independent games add via $\oplus$.
\end{itemize}
\end{theorem}
\begin{proof}[Proof for Nim]
Terminal $(0,\dots,0)$ has XOR sum $0$. Changing a single $n_i$ to $n_i' < n_i$ alters the highest bit where $n_i, n_i'$ differ, so XOR sum $0$ always moves to XOR sum $s \ne 0$. Conversely, if $s = n_1 \oplus \dots \oplus n_k \ne 0$, pick $n_i$ having a $1$ in the highest bit of $s$; then $n_i' = n_i \oplus s < n_i$, and replacing $n_i \to n_i'$ gives XOR sum $0$.
\end{proof}

\begin{problem}[Strategy Stealing in Chomp]
In Chomp on an $m \times n$ chocolate bar ($mn > 1$), players take turns choosing a remaining square $(i,j)$ and eating all $(i',j')$ with $i'\ge i, j'\ge j$; eating $(1,1)$ loses. Prove Player 1 has a winning strategy.
\end{problem}
\begin{solution}
Suppose Player 2 had a winning strategy. Let Player 1 eat only the corner $(m,n)$ on turn 1. Player 2's winning reply eats from some $(a,b)$. But eating $(a,b)$ on turn 1 would have produced the exact same board! Thus Player 1 can play $(a,b)$ on turn 1 and steal Player 2's winning strategy, contradiction.
\end{solution}

\section{Probability and the Probabilistic Method [Useful]}

\begin{definition}
A finite probability space assigns weights $p(\omega)\ge0$, $\sum p(\omega)=1$,
to outcomes. For an event $A$, $P(A) = \sum_{\omega\in A}p(\omega)$. A random
variable $X$ has \textbf{expected value} $E[X] = \sum_\omega p(\omega)X(\omega)$.
\end{definition}

\begin{theorem}[Linearity of Expectation, Alteration, and Lovász Local Lemma]
\leavevmode
\begin{enumerate}[label=(\alph*)]
\item \textbf{(Linearity)} For any random variables (independent or not!), $E[X + Y] = E[X] + E[Y]$ and $E[cX] = cE[X]$.
\item \textbf{(Existence \& Union Bound)} Some outcome has $X(\omega) \ge E[X]$ and some has $X(\omega) \le E[X]$. If $\sum P(A_i) < 1$, then $P(\bigcap \overline{A_i}) > 0$.
\item \textbf{(Alteration / Deletion Method)} If a random sample $S$ has expected size $E[X]$ and $E[Y]$ bad substructures, deleting one element per bad substructure leaves a clean subset of size $\ge E[X] - E[Y]$.
\item \textbf{(Lovász Local Lemma)} If bad events $A_i$ satisfy $P(A_i) \le p$, each $A_i$ is mutually independent of all but at most $d$ other $A_j$, and $ep(d+1) \le 1$, then $P(\bigcap \overline{A_i}) > 0$.
\end{enumerate}
\end{theorem}

\begin{example}
A random permutation of $[n]$ has, on average, $1$ fixed point ($E[\sum X_i] = n \cdot \frac{1}{n} = 1$).
\end{example}

\begin{problem}[Caro--Wei Theorem and Alteration Bound]
Let $G = (V,E)$ be a graph with $n$ vertices, $m$ edges, and average degree $d = 2m/n \ge 1$.
(a) Prove $G$ has an independent set of size at least $\sum_{v \in V}\frac{1}{\deg(v)+1}$.
(b) Use the alteration method to prove $G$ has an independent set of size at least $\frac{n}{2d}$.
\end{problem}
\begin{solution}
(a) Pick a uniformly random permutation $\pi$ of $V$ and let $I$ be the set of vertices $v$ that appear in $\pi$ before all of their neighbors. No edge can have both endpoints in $I$, so $I$ is independent. Each $v$ is first among $\{v\}\cup N(v)$ with probability $\frac{1}{\deg(v)+1}$, so $E[|I|] = \sum_{v\in V}\frac{1}{\deg(v)+1}$.
(b) Pick each vertex independently with probability $p = 1/d$, and delete one endpoint from each edge with both endpoints chosen. The expected remaining size is $np - mp^2 = np - \frac{nd}{2}p^2 = \frac{n}{2d}$.
\end{solution}

\section{Graph Theory}

\begin{definition}
A \textbf{graph} $G=(V,E)$ consists of vertices and edges. The \textbf{degree} $\deg(v)$ is the number of edges at $v$. A \textbf{tree} is a connected graph with no cycles.
\end{definition}

\begin{theorem}[Handshake Lemma and Trees]
$\sum_{v\in V}\deg(v) = 2|E|$. A tree on $n$ vertices has $n-1$ edges; every connected graph on $n$ vertices has at least $n-1$ edges.
\end{theorem}
\begin{proof}
Each edge contributes $2$ to $\sum\deg(v)$. A tree has a leaf (endpoint of a maximal path); remove it and induct.
\end{proof}

\subsection{Eulerian, Hamiltonian, and Planar Graphs}

\begin{theorem}[Eulerian Circuits, Dirac's Theorem, and Euler's Planar Formula]
\leavevmode
\begin{enumerate}[label=(\alph*)]
\item \textbf{(Euler)} A connected graph has an Eulerian circuit iff every vertex has even degree.
\item \textbf{(Dirac)} If $n \ge 3$ and $\deg(v) \ge n/2$ for all $v$, then $G$ has a Hamiltonian cycle.
\item \textbf{(Euler's Planar Formula)} Any connected plane graph satisfies $V - E + F = 2$. For $V \ge 3$, $3F \le 2E \implies E \le 3V - 6$ ($E \le 2V - 4$ if triangle-free), so every planar graph has a vertex of degree $\le 5$.
\end{enumerate}
\end{theorem}
\begin{proof}[Proof of (c)]
For a tree, $E=V-1$ and $F=1$, so $V-E+F=2$. Deleting an edge on a cycle reduces $E$ and $F$ by $1$ each, preserving $V-E+F$.
\end{proof}

\begin{problem}[Hamiltonian Path in a Tournament]
Prove that every tournament on $n$ players (every pair plays, one wins) contains a Hamiltonian path $p_1 \to p_2 \to \dots \to p_n$.
\end{problem}
\begin{solution}
Induct on $n$. Given $p_1 \to \dots \to p_{n-1}$ and a new player $v$: if $v$ beats $p_1$, prepend $v$; if $p_{n-1}$ beats $v$, append $v$. Otherwise let $i$ be the first index with $v$ beating $p_{i+1}$; then $p_i$ beats $v$, so insert $v$ between $p_i$ and $p_{i+1}$.
\end{solution}

\subsection{Hall's Marriage Theorem, König's Theorem, and Stable Matching [Advanced]}

\begin{theorem}[Hall, König, and Gale--Shapley]
\leavevmode
\begin{enumerate}[label=(\alph*)]
\item \textbf{(Hall's Marriage Theorem)} A bipartite graph $(A, B)$ has a matching covering $A$ iff $|N(S)| \ge |S|$ for every $S \subseteq A$.
\item \textbf{(König)} In any bipartite graph, max matching size $=$ min vertex cover size.
\item \textbf{(Gale--Shapley)} Given $n$ men and $n$ women with strict preferences, the propose-and-reject algorithm always terminates in a stable matching.
\end{enumerate}
\end{theorem}
\begin{proof}[Proof of Hall]
Necessity is clear. For sufficiency, induct on $|A|$. If $|N(S)| \ge |S|+1$ for all nonempty proper $S \subsetneq A$, match any edge $ab$ and induct on $A\setminus\{a\}$. Otherwise some proper $S_0$ has $|N(S_0)| = |S_0|$. Match $S_0$ onto $N(S_0)$ by induction; for $T \subseteq A\setminus S_0$, $|N(T)\setminus N(S_0)| = |N(T\cup S_0)| - |N(S_0)| \ge |T\cup S_0| - |S_0| = |T|$, so Hall's condition holds on $(A\setminus S_0, B\setminus N(S_0))$.
\end{proof}

\begin{problem}
A deck of $mn$ cards ($n$ values, each repeated $m$ times) is dealt into $n$ piles of $m$ cards. Prove one can pick one card from each pile so that all $n$ values appear.
\end{problem}
\begin{solution}
Form a bipartite graph between piles and values. Any $k$ piles hold $km$ cards; since each value appears only $m$ times in total, at least $\lceil km/m \rceil = k$ distinct values appear in those $k$ piles. By Hall's theorem, a perfect matching exists.
\end{solution}

\subsection{Extremal Graph Theory and Ramsey Theory [Advanced]}

\begin{theorem}[Mantel, Turán, Kővári--Sós--Turán, and Ramsey]
\leavevmode
\begin{enumerate}[label=(\alph*)]
\item \textbf{(Mantel)} A triangle-free graph on $n$ vertices has at most $\lfloor n^2/4 \rfloor$ edges.
\item \textbf{(Turán)} A $K_{r+1}$-free graph on $n$ vertices has at most $\bigl(1 - \frac1r\bigr)\frac{n^2}{2}$ edges (achieved by the balanced $r$-partite Turán graph).
\item \textbf{(Kővári--Sós--Turán)} A $C_4$-free graph on $n$ vertices satisfies $|E| \le \frac{n}{4}(1+\sqrt{4n-3}) = O(n^{3/2})$.
\item \textbf{(Ramsey)} Every $2$-coloring of $K_N$ with $N \ge R(s,t)$ has a red $K_s$ or blue $K_t$, where $R(s,t) \le R(s-1,t) + R(s,t-1) \le \binom{s+t-2}{s-1}$ and $R(3,3) = 6$.
\end{enumerate}
\end{theorem}
\begin{proof}
(a) For each edge $uv$, $\deg u + \deg v \le n$; summing over edges gives $\sum_{v} (\deg v)^2 = \sum_{uv}(\deg u + \deg v) \le n|E|$. By Cauchy--Schwarz, $\frac{4|E|^2}{n} \le \sum (\deg v)^2 \le n|E| \implies |E| \le n^2/4$.
(c) Count $2$-edge paths $(v, \{u,w\})$: since $G$ is $C_4$-free, each $\{u,w\}$ has at most $1$ common neighbor, so $n\binom{2|E|/n}{2} \le \sum_v \binom{\deg v}{2} \le \binom{n}{2}$.
(d) A vertex in $K_{R(s-1,t)+R(s,t-1)}$ has either $\ge R(s-1,t)$ red neighbors or $\ge R(s,t-1)$ blue neighbors; induct.
\end{proof}

\section{Algebraic Combinatorics and Combinatorial Geometry [Advanced]}

\subsection{Linear Algebra (Oddtown), Nullstellensatz, and Burnside's Lemma}

\begin{theorem}[Oddtown, Combinatorial Nullstellensatz, and Burnside's Lemma]
\leavevmode
\begin{enumerate}[label=(\alph*)]
\item \textbf{(Oddtown)} If $A_1,\dots,A_m \subseteq [n]$ have $|A_i|$ odd and $|A_i\cap A_j|$ even ($i\ne j$), then $m \le n$. \emph{Proof:} Incidence vectors $v_i \in \FF_2^n$ satisfy $v_i \cdot v_j = \delta_{ij} \pmod 2$, so they are linearly independent in $\FF_2^n$.
\item \textbf{(Alon's Combinatorial Nullstellensatz)} If $\deg P = \sum d_i$ and $[x_1^{d_1}\cdots x_n^{d_n}]P \ne 0$, then for any $S_i \subseteq F$ with $|S_i| > d_i$, there exists $s \in S_1\times\dots\times S_n$ with $P(s) \ne 0$ (reduce $P$ modulo $\prod_{s\in S_i}(x_i-s)$; the leading coefficient survives so the remainder is nonzero).
\item \textbf{(Cauchy--Davenport)} For nonempty $A, B \subseteq \ZZ/p\ZZ$, $|A+B| \ge \min(p, |A|+|B|-1)$ (apply Nullstellensatz to $\prod_{c\in C}(x+y-c)$ where $[x^{|A|-1}y^{|B|-1}] = \binom{|A|+|B|-2}{|A|-1} \not\equiv 0 \pmod p$).
\item \textbf{(Burnside's Lemma)} The number of orbits of a finite group $G$ acting on $X$ is $|X/G| = \frac{1}{|G|}\sum_{g\in G}|X^g|$ (double count $\{(g,x) : g\cdot x = x\}$).
\end{enumerate}
\end{theorem}

\subsection{Convex Hulls, Pick's Theorem, and Helly's Theorem}

\begin{theorem}[Convex Hulls, Pick's Theorem, and Helly's Theorem]
\leavevmode
\begin{enumerate}[label=(\alph*)]
\item For finite $S \subset \RR^2$, the \textbf{convex hull} is the set of convex combinations $\sum \lambda_i p_i$ ($\lambda_i \ge 0, \sum \lambda_i = 1$); its vertices lie in $S$.
\item \textbf{(Pick's Theorem)} The area of a simple lattice polygon is $A = I + \frac{B}{2} - 1$.
\item \textbf{(Helly's Theorem)} If every $d+1$ sets in a finite family of convex sets in $\RR^d$ have a common point, then all sets in the family have a common point.
\end{enumerate}
\end{theorem}

\begin{problem}
Given $2n$ points in general position (no 3 collinear), $n$ red and $n$ blue,
prove they can be split into $n$ red--blue pairs whose connecting segments are
pairwise non-crossing.
\end{problem}
\begin{solution}
Among all perfect red--blue pairings, choose one minimizing the total length of
the $n$ segments. If two segments $R_1B_1$ and $R_2B_2$ crossed at $X$, then by the triangle inequality
$R_1B_2 + R_2B_1 < (R_1X + XB_2) + (R_2X + XB_1) = R_1B_1 + R_2B_2$,
so swapping partners strictly decreases total length --- contradiction.
\end{solution}
%======================================================================
\part{Geometry}
%======================================================================

Throughout, $ABC$ is a triangle with side lengths $a=BC$, $b=CA$, $c=AB$,
semiperimeter $s = \frac{a+b+c}{2}$, circumradius $R$, inradius $r$, and area $[ABC]$.

\section{Triangle Centers, the Nine-Point Circle, and Fact 5}

\begin{definition}\leavevmode
\begin{itemize}
\item \textbf{Centroid $G$}: intersection of medians; divides each median $2:1$
from the vertex.
\item \textbf{Incenter $I$}: intersection of internal angle bisectors; center of
the inscribed circle, equidistant from the sides.
\item \textbf{Circumcenter $O$}: intersection of perpendicular bisectors;
center of the circle through $A,B,C$.
\item \textbf{Orthocenter $H$}: intersection of the altitudes.
\item \textbf{Excenter $I_A$}: intersection of the internal bisector from $A$
and external bisectors from $B, C$; center of the excircle opposite $A$.
\end{itemize}
\end{definition}

\begin{theorem}[Concurrency of Medians]
The three medians meet at a point $G$ dividing each in ratio $2:1$.
\end{theorem}
\begin{proof}
With position vectors, let $G = \frac{A+B+C}{3}$. The midpoint of $BC$ is
$M = \frac{B+C}{2}$, and $G = A + \frac{2}{3}(M - A)$ lies on median $AM$ at
ratio $2:1$; by symmetry $G$ lies on all three medians.
\end{proof}

\begin{theorem}[Euler Line]
$O$, $G$, $H$ are collinear with $\vec{OH} = 3\,\vec{OG}$, equivalently
$H = A + B + C$ when $O$ is the origin.
\end{theorem}
\begin{proof}
Take $O$ as origin, so $|A|=|B|=|C|=R$. Let $H' = A+B+C$. Then
$H' - A = B + C$ is perpendicular to $B - C$ since
$(B+C)\cdot(B-C) = |B|^2 - |C|^2 = 0$; thus $AH' \perp BC$, i.e.\ $H'$ lies on
the altitude from $A$, and by symmetry on all altitudes, so $H'=H$. Since
$G = \frac{A+B+C}{3} = \frac{H}{3}$, the collinearity and ratio follow.
\end{proof}

\begin{theorem}[Nine-Point Circle]
The midpoints of the sides, the feet of the altitudes, and the midpoints of
$AH$, $BH$, $CH$ all lie on one circle with center $N$ (the midpoint of $OH$)
and radius $R/2$.
\end{theorem}
\begin{proof}
Consider the homothety $h$ centered at $H$ with ratio $\frac12$. It maps the
circumcircle to a circle of radius $R/2$ centered at the midpoint $N$ of $OH$.
It maps $A, B, C$ to the midpoints of $HA, HB, HC$: three of the nine points.
Next, the reflection of $H$ over the midpoint $M_A$ of $BC$ is the antipode
$A'$ of $A$ on the circumcircle (indeed $BHCA'$ is a parallelogram: $BH \parallel CA'$
as both are perpendicular to $AC$, and $CH \parallel BA'$ likewise); hence
$h(A') = M_A$, putting the three side midpoints on the circle. Finally the
reflection of $H$ over the foot of the altitude from $A$ lies on the
circumcircle (reflecting $H$ across $BC$ gives a point $H_A$ with
$\angle BH_AC = \angle BHC = 180^\circ - \angle A$, so $H_A$ is concyclic with
$A,B,C$), hence $h(H_A)$ is the foot itself, and the three feet lie on the
circle too.
\end{proof}

\subsection{The Incenter--Excenter Lemma and Intouch/Extouch Symmetry}

\begin{theorem}[The Incenter--Excenter Lemma (``Fact 5'')]
Let $M$ be the midpoint of arc $BC$ of $(ABC)$ not containing $A$. Then $M$ is the center of a circle passing through $B, I, C, I_A$:
\[
MB = MC = MI = MI_A.
\]
\end{theorem}
\begin{proof}
Since $M$ bisects arc $BC$, line $AM$ bisects $\angle A$ (so $A, I, M, I_A$ are collinear) and $MB = MC$. By angle chasing,
$\angle MBI = \angle MBC + \angle CBI = \angle MAC + \frac{\angle B}{2} = \frac{\angle A + \angle B}{2}$,
while by the exterior angle theorem on $\triangle ABI$, $\angle BIM = \frac{\angle A + \angle B}{2}$. Thus $MB = MI$. Furthermore, internal and external bisectors at $B$ and $C$ are perpendicular, so $\angle IBI_A = \angle ICI_A = 90^\circ$. Hence $B, I, C, I_A$ lie on the circle with diameter $II_A$, whose center on $II_A$ must be $M$ since $MB = MI = MC$.
\end{proof}

\begin{proposition}[Tangent Lengths and Euler's Formula $OI^2 = R(R-2r)$]
\leavevmode
\begin{enumerate}[label=(\alph*)]
\item If the incircle touches $BC, CA, AB$ at $D, E, F$ and the $A$-excircle touches $BC$ at $D_a$, then $AE=AF=s-a$, $BD=BF=s-b$, $CD=CE=s-c$, and $BD = CD_a = s-b$ (so $D$ and $D_a$ are symmetric about the midpoint of $BC$).
\item $OI^2 = R^2 - 2Rr$, so $R \ge 2r$.
\end{enumerate}
\end{proposition}
\begin{proof}[Proof of (b)]
Let $AI$ meet $(ABC)$ at $M$, and let $M'$ be the antipode of $M$ on $(ABC)$ and $D$ the foot from $I$ to $AB$ ($ID=r$). By power of a point and Fact 5, $R^2 - OI^2 = AI \cdot IM = AI \cdot BM$. Since $\triangle AID \sim \triangle M'MB$, $\frac{AI}{2R} = \frac{r}{BM} \implies AI \cdot BM = 2Rr$.
\end{proof}

\begin{problem}
Prove that the reflection of the orthocenter $H$ over side $BC$ lies on the
circumcircle of $ABC$.
\end{problem}
\begin{solution}
In the quadrilateral formed by the altitudes from $B$ and $C$, $\angle BHC = 180^\circ - \angle A$.
Reflecting $H$ over $BC$ to $H'$ preserves the angle, so
$\angle BH'C = \angle BHC = 180^\circ - \angle BAC$, meaning $ABH'C$ is cyclic.
\end{solution}

\section{Angle Chasing, Directed Angles, and Cyclic Quadrilaterals}

\begin{theorem}[Inscribed Angle Theorem]
An inscribed angle equals half the central angle subtending the same arc.
Consequently: angles subtending the same arc are equal, and an angle in a
semicircle is right.
\end{theorem}
\begin{proof}
If one side of the inscribed angle at $P$ is a diameter $PQ'$, $\triangle OPQ$ is isosceles ($OP=OQ=R$), so the central angle $\angle Q'OQ = 2\angle OPQ$. The general case adds or subtracts two such angles.
\end{proof}

\textbf{Directed Angles Modulo $180^\circ$.} Write $\dang(\ell_1, \ell_2)$ for the counterclockwise angle from line $\ell_1$ to line $\ell_2$ modulo $180^\circ$, and $\dang ABC := \dang(BA, BC)$. Then $\dang ABC = -\dang CBA$ and $\dang AXP + \dang PXB = \dang AXB$.

\begin{theorem}[Cyclic Quadrilateral Criteria]
For a convex quadrilateral $ABCD$, the following are equivalent:
(i) $ABCD$ is cyclic; (ii) $\angle ABD = \angle ACD$; (iii) $\angle ABC + \angle CDA = 180^\circ$; (iv) for any four non-collinear points, $\dang ABD = \dang ACD$.
\end{theorem}
\begin{proof}
(i)$\Rightarrow$(ii)--(iv) by the inscribed angle theorem. Conversely, if $\dang ABD = \dang ACD$ and line $AC$ meets circle $(ABD)$ again at $C'$, then $\dang AC'D = \dang ABD = \dang ACD$, so lines $C'D$ and $CD$ are parallel and share $D$, forcing $C' = C$.
\end{proof}

\begin{theorem}[Ptolemy's Theorem]
For a cyclic quadrilateral $ABCD$: $AC \cdot BD = AB\cdot CD + AD \cdot BC$.
In general (Ptolemy's inequality) $AC\cdot BD \le AB\cdot CD + AD\cdot BC$
with equality iff $ABCD$ is cyclic in this order.
\end{theorem}
\begin{proof}
Choose $E$ on diagonal $AC$ with $\angle ABE = \angle DBC$. Since
$\angle BAE = \angle BDC$, triangles $ABE$ and
$DBC$ are similar, so $AE\cdot BD = AB\cdot CD$. Also $\angle EBC = \angle ABD$ and
$\angle BCE = \angle BDA$, so $\triangle EBC \sim \triangle ABD$, giving $EC \cdot BD = AD \cdot BC$.
Adding gives $(AE+EC)\cdot BD = AC\cdot BD = AB\cdot CD + AD\cdot BC$.
\end{proof}

\begin{theorem}[Simson Line and Miquel's Theorem]
\leavevmode
\begin{enumerate}[label=(\alph*)]
\item \textbf{(Simson Line)} The feet $X, Y, Z$ of the perpendiculars from $P$ to $BC, CA, AB$ are collinear iff $P$ lies on $(ABC)$.
\item \textbf{(Miquel's Theorem)} If $D, E, F$ lie on lines $BC, CA, AB$, the circles $(AEF)$, $(BFD)$, and $(CDE)$ concur at a point $M$.
\end{enumerate}
\end{theorem}
\begin{proof}
(a) Since $P,Y,Z,A$ and $P,X,Z,B$ are concyclic, $\dang YZX = \dang YZP + \dang PZX = \dang YAP + \dang PBX = \dang CAP - \dang CBP$, which is $0 \pmod{180^\circ}$ iff $P \in (ABC)$.
(b) If $(AEF)$ and $(BFD)$ meet at $M \ne F$, then $\dang DME = \dang DMF + \dang FME = \dang DBF + \dang FAE = \dang DCE$, so $M \in (CDE)$.
\end{proof}

\begin{problem}
Point $P$ lies on the arc $BC$ of the circumcircle of equilateral triangle
$ABC$ (not containing $A$). Prove $PA = PB + PC$.
\end{problem}
\begin{solution}
Apply Ptolemy to cyclic quadrilateral $ABPC$:
$PA \cdot BC = PB \cdot AC + PC \cdot AB$. Since $AB=BC=CA$, divide through to
get $PA = PB + PC$.
\end{solution}

\section{Similar Triangles and Metric Relations}

\begin{theorem}[Similarity Criteria and Stewart's Theorem]
\leavevmode
\begin{enumerate}[label=(\alph*)]
\item Triangles are similar under \textbf{AA}, \textbf{SAS}, or \textbf{SSS}; areas scale as the square of the similarity ratio.
\item \textbf{(Stewart's Theorem)} For $D \in BC$ with $BD=m, DC=n, m+n=a$, and $AD=d$:
\[
man + dad = bmb + cnc \quad \bigl(a(mn + d^2) = b^2 m + c^2 n\bigr).
\]
Consequently, median $m_a^2 = \frac{2b^2+2c^2-a^2}{4}$ and internal bisector $\ell_a^2 = bc\bigl(1 - \frac{a^2}{(b+c)^2}\bigr)$.
\item \textbf{(Heron \& Brahmagupta)} $[ABC] = \sqrt{s(s-a)(s-b)(s-c)}$, and a cyclic quad with sides $a,b,c,d$ has area $\sqrt{(s-a)(s-b)(s-c)(s-d)}$.
\end{enumerate}
\end{theorem}
\begin{proof}[Proof of Stewart]
Let $\theta = \angle ADB$; by the Law of Cosines on $\triangle ABD$ and $\triangle ACD$, $c^2 = m^2+d^2-2md\cos\theta$ and $b^2 = n^2+d^2+2nd\cos\theta$. Multiply the first by $n$, the second by $m$, and add.
\end{proof}

\begin{example}
In $\triangle ABC$, $DE \parallel BC$ with $D\in AB, E\in AC$. If $AD=3, DB=2, BC=10$, then $\triangle ADE \sim \triangle ABC$ with ratio $\frac{3}{5}$, so $DE = 6$.
\end{example}

\begin{problem}
In right triangle $ABC$ ($\angle C = 90^\circ$) with altitude $CH \perp AB$, prove $CH^2 = AH \cdot BH$ and $AC^2 = AH \cdot AB$.
\end{problem}
\begin{solution}
Since $\angle ACH = 90^\circ - \angle A = \angle B$, $\triangle ACH \sim \triangle CBH \sim \triangle ABC$ (AA), giving $\frac{CH}{BH}=\frac{AH}{CH}$ and $\frac{AC}{AB}=\frac{AH}{AC}$.
\end{solution}

\section{Power of a Point, Radical Axes, and Apollonius Circles}

\begin{definition}
The \textbf{power} of point $P$ with respect to circle $\omega(O,r)$ is $\pow(P,\omega) = OP^2 - r^2$.
\end{definition}

\begin{theorem}[Power of a Point, Radical Center, and Circle of Apollonius]
\leavevmode
\begin{enumerate}[label=(\alph*)]
\item If a line through $P$ meets $\omega$ at $X, Y$, then $\overrightarrow{PX}\cdot\overrightarrow{PY} = OP^2 - r^2$ (and $= PT^2$ for tangent $PT$). Conversely, if $\overrightarrow{PX}\cdot\overrightarrow{PY} = \overrightarrow{PZ}\cdot\overrightarrow{PW}$ on two lines through $P$, then $X,Y,Z,W$ are concyclic.
\item \textbf{(Radical Axis \& Coaxial Pencils)} The locus of points with equal power to non-concentric $\omega_1, \omega_2$ is a line $\perp O_1O_2$ (the \textbf{radical axis}). For three circles, the pairwise radical axes concur at the \textbf{radical center}. A disjoint pencil $\omega_1 + \lambda\omega_2 = 0$ has two radius-$0$ limiting points $L_1, L_2$; inverting at $L_1$ maps the pencil to concentric circles.
\item \textbf{(Circle of Apollonius)} For distinct $A, B$ and $k > 0$ ($k \ne 1$), the locus $\{P : PA/PB = k\}$ is the circle with diameter $DD'$, where $D, D' \in AB$ are the internal and external angle bisector feet ($\frac{DA}{DB} = \frac{D'A}{D'B} = k$), since $PD, PD'$ bisect $\angle APB \iff \angle DPD' = 90^\circ$.
\end{enumerate}
\end{theorem}
\begin{proof}[Proof of (a)]
With $M$ the midpoint of $XY$, $OM \perp XY$ and $PX\cdot PY = PM^2 - MX^2 = (OP^2 - OM^2) - (r^2 - OM^2) = OP^2 - r^2$.
\end{proof}

\begin{problem}
Two circles intersect at $A$ and $B$. A line through $B$ meets the circles
again at $C$ and $D$. Tangents to the circles at $C$ and $D$ meet at $P$.
Prove that $P$, $A$, $C$, $D$ are concyclic.
\end{problem}
\begin{solution}
By tangent--chord angles, $\angle PCB = \angle CAB$ and $\angle PDB = \angle DAB$. Thus
$\angle CPD = 180^\circ - \angle PCD - \angle PDC = 180^\circ - \angle CAB - \angle DAB = 180^\circ - \angle CAD$,
so $PCAD$ is cyclic.
\end{solution}

\section{Homothety, Spiral Similarity, and Complete Quadrilaterals}

\begin{definition}
A \textbf{homothety} $h(O, k)$ maps $P \mapsto O + k(P - O)$. A \textbf{spiral similarity} centered at $O$ is a rotation composed with a homothety at $O$: $z \mapsto \alpha(z-o)+o$.
\end{definition}

\begin{theorem}[Monge, Spiral Similarity Lemma, and Complete Quadrilaterals]
\leavevmode
\begin{enumerate}[label=(\alph*)]
\item \textbf{(Monge's Theorem)} For three circles of distinct radii, the three external homothety centers are collinear.
\item \textbf{(Spiral Similarity Lemma)} If lines $AC$ and $BD$ meet at $X$, the second intersection $P$ of circles $(XAB)$ and $(XCD)$ is the center of the spiral similarity sending $A \mapsto B, C \mapsto D$ \emph{and} the spiral similarity sending $A \mapsto C, B \mapsto D$: $\triangle PAB \sim \triangle PCD$ and $\triangle PAC \sim \triangle PBD$.
\item \textbf{(Complete Quadrilateral)} Given four lines forming four triangles, their four circumcircles concur at the Miquel point $M$ (the spiral center), $M$ and the four circumcenters are concyclic, $M$ has a common Simson line, the four orthocenters are collinear, and the midpoints of the three diagonals are collinear (\textbf{Newton--Gauss line}).
\end{enumerate}
\end{theorem}
\begin{proof}[Proof of (b)]
In cyclic quads $PXAB$ and $PXCD$: $\dang PAB = \dang PXB = \dang PXD = \dang PCD$ and $\dang PBA = \dang PXA = \dang PXC = \dang PDC$, so $\triangle PAB \sim \triangle PCD$, which also implies $\triangle PAC \sim \triangle PBD$.
\end{proof}

\begin{problem}
Circles $\omega_1$ and $\omega_2$ are internally tangent at $T$, with
$\omega_1$ inside. A chord $AB$ of $\omega_2$ is tangent to $\omega_1$ at $P$.
Prove that line $TP$ bisects the arc $AB$ of $\omega_2$ not containing $T$.
\end{problem}
\begin{solution}
The homothety $h$ at $T$ taking $\omega_1 \to \omega_2$ maps $P$ to $M = h(P) \in \omega_2$, and maps tangent $AB$ of $\omega_1$ to the tangent of $\omega_2$ at $M$, which is parallel to $AB$. Thus $M$ is the midpoint of arc $AB$.
\end{solution}

\section{Ceva, Menelaus, Isogonal Conjugates, and Symmedians}

\begin{theorem}[Ceva and Menelaus]
For $D, E, F$ on lines $BC, CA, AB$:
\begin{itemize}
\item \textbf{(Ceva)} $AD, BE, CF$ concur iff $\dfrac{BD}{DC}\cdot\dfrac{CE}{EA}\cdot\dfrac{AF}{FB} = 1$ (signed ratios; proved via $\frac{BD}{DC} = \frac{[ABP]}{[ACP]}$).
\item \textbf{(Menelaus)} $D, E, F$ are collinear iff $\dfrac{BD}{DC}\cdot\dfrac{CE}{EA}\cdot\dfrac{AF}{FB} = -1$ (signed ratios; proved by dropping perpendiculars to $DEF$).
\end{itemize}
\end{theorem}

\begin{theorem}[Isogonal Conjugates and Symmedians]
\leavevmode
\begin{enumerate}[label=(\alph*)]
\item Reflecting concurrent cevians $AP, BP, CP$ across the angle bisectors yields concurrent cevians meeting at the \textbf{isogonal conjugate} $P^*$ (by Trig Ceva; e.g.\ $O^* = H$).
\item The isogonal conjugate of median $AM_A$ is the \textbf{$A$-symmedian} (concurring at $K = G^*$). If the tangents to $(ABC)$ at $B, C$ meet at $X$ and $AX$ meets $BC$ at $D$ and $(ABC)$ at $T$, then $AD$ is the $A$-symmedian, $\frac{BD}{DC} = \frac{c^2}{b^2}$, and $ABTC$ is a \textbf{harmonic quadrilateral} ($AB \cdot CT = AC \cdot BT$).
\end{enumerate}
\end{theorem}
\begin{proof}[Proof of (b)]
Using $XB=XC$ and $\angle ABX = 180^\circ - C, \angle ACX = 180^\circ - B$, the Law of Sines on $\triangle ABX$ and $\triangle ACX$ gives $\frac{\sin\angle BAX}{\sin\angle CAX} = \frac{\sin C}{\sin B} = \frac{c}{b} = \frac{\sin\angle M_AAC}{\sin\angle BAM_A}$, so $AX$ is the $A$-symmedian.
\end{proof}

\begin{problem}
Prove using Ceva that the internal angle bisectors of a triangle are concurrent.
\end{problem}
\begin{solution}
By the Angle Bisector Theorem, $\frac{BD}{DC} = \frac{c}{b}$, $\frac{CE}{EA} = \frac{a}{c}$, $\frac{AF}{FB} = \frac{b}{a}$; their product is $1$.
\end{solution}

\section{Trigonometric and Coordinate Methods [Useful]}

\begin{theorem}[Law of Sines, Law of Cosines, and Trig Ceva]
$\frac{a}{\sin A} = \frac{b}{\sin B} = \frac{c}{\sin C} = 2R$; $a^2 = b^2+c^2-2bc\cos A$; $[ABC] = \frac12 bc\sin A = \frac{abc}{4R} = rs$. Cevians $AD, BE, CF$ concur iff $\frac{\sin\angle BAD}{\sin\angle DAC}\frac{\sin\angle CBE}{\sin\angle EBA}\frac{\sin\angle ACF}{\sin\angle FCB} = 1$.
\end{theorem}

\begin{itemize}
\item \textbf{Cartesian \& Shoelace:} $[P_1\cdots P_n] = \frac12\bigl|\sum_i (x_i y_{i+1} - x_{i+1}y_i)\bigr|$.
\item \textbf{Complex Numbers on Unit Circle ($|a|=|b|=|c|=1$, $\bar a = 1/a$):} $o=0$, $g=\frac{a+b+c}{3}$, $h=a+b+c$; chord $ab$: $z+ab\bar z = a+b$; tangent at $a$: $z+a^2\bar z = 2a$; foot from $p$ to $ab$: $\frac12(a+b+p-ab\bar p)$.
\item \textbf{Barycentric Coordinates $P=(x:y:z)$:} $A=(1:0:0)$, $G=(1:1:1)$, $I=(a:b:c)$, $K=(a^2:b^2:c^2)$; circumcircle is $a^2yz + b^2zx + c^2xy = 0$.
\end{itemize}

\begin{problem}
With circumcenter $O=0$ and $|a|=|b|=|c|=1$, prove the reflection of $O$ across $AB$ is $a+b$, and that it is the circumcenter of $\triangle HAB$.
\end{problem}
\begin{solution}
The midpoint of $AB$ is $\frac{a+b}{2}$, so reflecting $0$ gives $O' = a+b$. With $h=a+b+c$, we have $|O'-a|=|b|=1$, $|O'-b|=|a|=1$, and $|O'-h|=|-c|=1$, so $O'$ is the circumcenter of $\triangle HAB$.
\end{solution}

\section{Inversion and Mixtilinear Incircles [Useful]}

\begin{theorem}[Inversion, Mixtilinear Incircles, and Casey's Theorem]
Under inversion at $O$ with radius $r$ ($OP\cdot OP^* = r^2$):
\begin{enumerate}[label=(\alph*)]
\item Lines through $O \leftrightarrow$ lines through $O$; circles through $O \leftrightarrow$ lines not through $O$; circles not through $O \leftrightarrow$ circles not through $O$; angles are preserved; and $P^*Q^* = \frac{r^2\cdot PQ}{OP\cdot OQ}$ ($\triangle OPQ \sim \triangle OQ^*P^*$).
\item \textbf{($\sqrt{bc}$-Inversion)} Inverting at $A$ with radius $\sqrt{bc}$ and reflecting across the bisector of $\angle A$ swaps $B\leftrightarrow C$ and swaps line $BC \leftrightarrow (ABC)$.
\item \textbf{(Mixtilinear Incircle)} If the $A$-mixtilinear incircle touches $AB, AC$ at $F, E$ and $(ABC)$ internally at $T$, then the incenter $I$ is the midpoint of $EF$ (by Pascal's theorem on $A B M_C T M_B C$, where $M_C, M_B$ are the arc midpoints of $AB, AC$).
\item \textbf{(Casey's Theorem)} Four circles are tangent to a common circle $\Omega$ in cyclic order iff their external bitangents satisfy $t_{12}t_{34} + t_{23}t_{14} = t_{13}t_{24}$.
\end{enumerate}
\end{theorem}

\begin{problem}
Three circles pass through $O$ and pairwise intersect again at $A, B, C$. Show that under inversion at $O$, $A^*, B^*, C^*$ form a triangle whose angles equal the pairwise intersection angles of the three circles at $O$.
\end{problem}
\begin{solution}
Inverting at $O$ maps the three circles through $O$ to three lines pairwise intersecting at $A^*, B^*, C^*$. Since inversion preserves angles between curves, the angles between the circles at $O$ equal the interior angles of $\triangle A^*B^*C^*$.
\end{solution}

\section{Projective Geometry [Useful]}

\begin{definition}
For collinear $A, B, C, D$, the \textbf{cross ratio} $(A,B;C,D) = \frac{\overline{CA}}{\overline{CB}} \cdot \frac{\overline{DB}}{\overline{DA}}$ is preserved by central projection; $(A,B;C,D) = -1$ is a \textbf{harmonic bundle}.
\end{definition}

\begin{theorem}[Harmonic Bundles, Poles and Polars, Brokard, Desargues, and Pascal]
\leavevmode
\begin{enumerate}[label=(\alph*)]
\item \textbf{(La Hire)} The \textbf{polar} of $P$ with respect to $\omega(O,r)$ is the line $\perp OP$ through $P^*$. $P$ lies on the polar of $Q$ iff $Q$ lies on the polar of $P$; if secant $PXY$ meets the polar of $P$ at $Q$, then $(X,Y;Q,P) = -1$.
\item \textbf{(Brokard's Theorem)} For cyclic $ABCD$ inscribed in $\omega(O,r)$, $P=AB\cap CD, Q=AD\cap BC, R=AC\cap BD$ form a self-polar triangle $\triangle PQR$, so $O$ is the orthocenter of $\triangle PQR$.
\item \textbf{(Desargues, Pappus, Pascal, \& Brianchon)} Two triangles are perspective from a point iff perspective from a line \textbf{(Desargues)}. Opposite sides of a hexagon inscribed in a conic (or pair of lines, \textbf{Pappus}) meet in three collinear points \textbf{(Pascal; proved on a circle via three applications of Menelaus and Power of a Point)}. Main diagonals of a circumscribed hexagon concur \textbf{(Brianchon)}.
\item \textbf{(Moving Points)} A degree-$d$ projective map between lines/conics is uniquely determined by $d+1$ test positions ($3$ positions for a standard projectivity $d=1$).
\end{enumerate}
\end{theorem}

\begin{problem}
In $\triangle ABC$ inscribed in $\omega$, the tangents at $B, C$ meet at $X$, the tangent at $A$ meets $BC$ at $T$, and $AX$ meets $BC$ at $K$. Prove $(B, C; K, T) = -1$ and $TA^2 = TB \cdot TC$.
\end{problem}
\begin{solution}
$T$ lies on the polar of $A$ (tangent at $A$) and the polar of $X$ (line $BC$), so by La Hire the polar of $T$ is line $AX$. Since the polar of $T$ meets secant $BC$ at $K$, $(B,C;K,T) = -1$. Also $\triangle TAB \sim \triangle TCA$ gives $TA^2 = TB\cdot TC$.
\end{solution}

\section{Conic Sections and 3D Geometry [Advanced]}

\begin{theorem}[Optical Reflection Property, Cuboid Embedding, and Dandelin Spheres]
\leavevmode
\begin{enumerate}[label=(\alph*)]
\item \textbf{(Ellipse Reflection)} The tangent at $P$ to an ellipse with foci $F_1, F_2$ ($PF_1+PF_2=2a$) makes equal angles with $PF_1$ and $PF_2$ (every $X \ne P$ on the tangent lies outside the ellipse, so $P$ minimizes $XF_1+XF_2$ along the tangent line).
\item \textbf{(3D Cuboid Embedding)} An \textbf{equifacial tetrahedron} $ABCD$ ($AB=CD=a, AC=BD=b, AD=BC=c$) inscribes in a rectangular box of dimensions $x,y,z$ ($x^2+y^2=a^2$, etc.) with circumradius $R = \frac{1}{2\sqrt{2}}\sqrt{a^2+b^2+c^2}$ and volume $V = xyz - 4(\frac{1}{6}xyz) = \frac{1}{3}xyz$.
\item \textbf{(Dandelin Spheres)} The two spheres tangent to a cone and to a secant plane $\Pi$ touch $\Pi$ at the foci $F_1, F_2$ of the conic section.
\end{enumerate}
\end{theorem}

\begin{problem}
Prove that the product of the distances from the two foci $(\pm c, 0)$ of the ellipse $\frac{x^2}{a^2}+\frac{y^2}{b^2}=1$ ($c^2=a^2-b^2$) to any tangent line equals $b^2$.
\end{problem}
\begin{solution}
At $P = (a\cos\theta, b\sin\theta)$, the tangent is $bx\cos\theta + ay\sin\theta - ab = 0$. The product of the distances from $(\pm c, 0)$ is
\[
d_1 d_2 = \frac{|a^2b^2 - b^2c^2\cos^2\theta|}{b^2\cos^2\theta + a^2\sin^2\theta}
= \frac{b^2(a^2\sin^2\theta + b^2\cos^2\theta)}{b^2\cos^2\theta + a^2\sin^2\theta} = b^2. \qedhere
\]
\end{solution}
%======================================================================
\part{Number Theory}
%======================================================================

\section{Divisibility, GCD, and the Euclidean Algorithm}

\begin{definition}
For integers $a, b$ not both zero, $\gcd(a,b)$ is the largest positive integer
dividing both. Integers are \textbf{coprime} if $\gcd(a,b)=1$.
\end{definition}

\begin{theorem}[Division Algorithm]
For integers $a$ and $b > 0$ there are unique $q, r$ with $a = qb + r$ and
$0 \le r < b$.
\end{theorem}

\begin{theorem}[Euclidean Algorithm]
$\gcd(a, b) = \gcd(b, a - b) = \gcd(b, a \bmod b)$. Iterating terminates and
computes the gcd.
\end{theorem}
\begin{proof}
Any common divisor of $a, b$ divides $a - qb$ and conversely any common
divisor of $b, a-qb$ divides $a$; so the two pairs have identical sets of
common divisors, hence the same gcd. The remainders strictly decrease, so the
process reaches $\gcd(g, 0) = g$.
\end{proof}

\begin{theorem}[Bézout's Identity]
For any $a, b$ not both zero there exist integers $x, y$ with
$ax + by = \gcd(a,b)$. Consequently $\gcd(a,b)$ is the least positive element
of $\{ax+by : x,y \in \ZZ\}$, and $\gcd(a,b)=1$ iff $ax+by=1$ is solvable.
\end{theorem}
\begin{proof}
Let $d$ be the least positive integer of the form $ax + by$. Dividing,
$a = qd + r$ with $0 \le r < d$; then $r = a - qd$ is again of the form
$ax'+by'$, so minimality forces $r = 0$, i.e.\ $d \mid a$; similarly
$d \mid b$. Any common divisor of $a, b$ divides $ax+by = d$, so
$d = \gcd(a,b)$.
\end{proof}

\begin{corollary}[Euclid's Lemma]
If $p$ is prime and $p \mid ab$, then $p \mid a$ or $p \mid b$. More generally
if $\gcd(c, a) = 1$ and $c \mid ab$ then $c \mid b$.
\end{corollary}
\begin{proof}
If $p \nmid a$ then $\gcd(p,a)=1$, so $px + ay = 1$; multiply by $b$:
$b = pbx + aby$, and $p$ divides both terms on the right.
\end{proof}

\begin{example}
Compute $\gcd(2026, 748)$.
\end{example}
\begin{solution}
$2026 = 2\cdot748 + 530$; $748 = 1\cdot530 + 218$; $530 = 2\cdot218 + 94$;
$218 = 2\cdot 94 + 30$; $94 = 3\cdot30 + 4$; $30 = 7\cdot4+2$; $4 = 2\cdot2$.
So $\gcd = 2$.
\end{solution}

\begin{problem}
Prove that $\gcd(2^m - 1,\, 2^n - 1) = 2^{\gcd(m,n)} - 1$.
\end{problem}
\begin{solution}
Assume $m \ge n$. Then
$2^m - 1 = 2^{m-n}(2^n - 1) + (2^{m-n} - 1)$, so
$\gcd(2^m - 1, 2^n - 1) = \gcd(2^n - 1, 2^{m-n}-1)$: the exponents perform the
Euclidean algorithm, terminating at
$\gcd(2^{\gcd(m,n)}-1,\, 2^0 - 1) = 2^{\gcd(m,n)} - 1$.
\end{solution}

\section{Primes, Unique Factorization, and Multiplicative Functions}

\begin{theorem}[Euclid]
There are infinitely many primes.
\end{theorem}
\begin{proof}
If $p_1,\dots,p_k$ were all the primes, then $N = p_1\cdots p_k + 1 > 1$ has a
prime factor, but $N \equiv 1 \pmod{p_i}$ for each $i$ --- contradiction.
\end{proof}

\begin{theorem}[Fundamental Theorem of Arithmetic]
Every integer $n \ge 2$ factors into primes, uniquely up to order.
\end{theorem}
\begin{proof}
Existence follows by strong induction. For uniqueness,
suppose $p_1\cdots p_r = q_1 \cdots q_s$ are two factorizations of the
smallest $n$ with distinct factorizations. Then $p_1 \mid q_1\cdots q_s$, so
by Euclid's Lemma $p_1 \mid q_j$ for some $j$, forcing $p_1 = q_j$.
Cancelling gives a smaller counterexample, contradiction.
\end{proof}

With $n = \prod p_i^{a_i}$ and $m = \prod p_i^{b_i}$:
\[
\gcd(m,n) = \prod p_i^{\min(a_i,b_i)}, \quad
\lcm(m,n) = \prod p_i^{\max(a_i,b_i)}, \quad
d(n) = \prod (a_i + 1), \quad
\sigma(n) = \prod \frac{p_i^{a_i+1}-1}{p_i - 1}.
\]

\subsection{Multiplicative Functions and Möbius Inversion}
A function $f\colon\NN\to\CC$ is \textbf{multiplicative} if $f(1)=1$ and $f(mn)=f(m)f(n)$ whenever $\gcd(m,n)=1$. Euler's totient satisfies $\varphi(n) = n\prod_{p\mid n}\bigl(1 - \frac{1}{p}\bigr)$ and $\sum_{d\mid n}\varphi(d) = n$.

\begin{definition}[Möbius Function]
$\mu(n) = (-1)^k$ if $n$ is a product of $k$ distinct primes ($k\ge 0$), and $\mu(n) = 0$ if $p^2 \mid n$ for some prime $p$. It satisfies $\sum_{d\mid n}\mu(d) = [n=1]$.
\end{definition}

\begin{theorem}[Möbius Inversion Formula]
$$g(n) = \sum_{d \mid n} f(d) \quad\forall n \ge 1 \iff f(n) = \sum_{d \mid n} \mu(d) g\left(\frac{n}{d}\right) \quad\forall n \ge 1$$
\end{theorem}
\begin{proof}
$\sum_{d\mid n}\mu(d)g(n/d) = \sum_{d\mid n}\mu(d)\sum_{c\mid (n/d)}f(c) = \sum_{c\mid n}f(c)\sum_{d\mid (n/c)}\mu(d) = f(n)$.
\end{proof}

\begin{problem}
(a) Prove that $\sqrt[k]{n}$ is irrational unless $n$ is a perfect $k$-th power.
(b) Prove that $\sum_{d \mid n}\varphi(d) = n$.
\end{problem}
\begin{solution}
(a) If $\sqrt[k]{n} = a/b$, then $nb^k = a^k$, so $\nu_p(n) + k\nu_p(b) = k\nu_p(a) \implies k \mid \nu_p(n)$ for every prime $p$.
(b) Group the $n$ fractions $\frac{1}{n}, \frac{2}{n}, \dots, \frac{n}{n}$ by their reduced denominator $d \mid n$ when written in lowest terms $\frac{a}{d}$ ($\gcd(a,d)=1$); each $d \mid n$ appears $\varphi(d)$ times.
\end{solution}

\section{Modular Arithmetic, Orders, and Primitive Roots}

\begin{theorem}[Fermat's Little Theorem, Euler, CRT, and Wilson]
\leavevmode
\begin{enumerate}[label=(\alph*)]
\item \textbf{(Fermat / Euler)} If $\gcd(a,n)=1$, then $a^{\varphi(n)} \equiv 1 \pmod n$; for prime $p \nmid a$, $a^{p-1} \equiv 1 \pmod p$ (since $x \mapsto ax$ permutes the invertible residues).
\item \textbf{(Chinese Remainder Theorem)} If $n_1,\dots,n_k$ are pairwise coprime, $x \equiv a_i \pmod{n_i}$ has a unique solution modulo $N = n_1\cdots n_k$, given by $x = \sum_{i=1}^k a_i N_i M_i$ where $N_i = N/n_i$ and $N_i M_i \equiv 1 \pmod{n_i}$.
\item \textbf{(Wilson)} $p > 1$ is prime iff $(p-1)! \equiv -1 \pmod p$ (pair $x \in \{2,\dots,p-2\}$ with $x^{-1}$).
\end{enumerate}
\end{theorem}

\subsection{Orders and Primitive Roots}

\begin{definition}
If $\gcd(a, n) = 1$, the \textbf{order} $\ord_n(a)$ is the least $d \ge 1$
with $a^d \equiv 1 \pmod n$. An integer $g$ with $\ord_n(g) = \varphi(n)$ is a \textbf{primitive root} modulo $n$.
\end{definition}

\begin{theorem}[Fundamental Property of Order and Existence of Primitive Roots]\label{thm:order}
\leavevmode
\begin{enumerate}[label=(\alph*)]
\item $a^m \equiv 1 \pmod n \iff \ord_n(a) \mid m$. In particular $\ord_n(a) \mid \varphi(n)$.
\item Primitive roots modulo $n$ exist if and only if $n = 1, 2, 4, p^k,$ or $2p^k$ for an odd prime $p$. Every prime $p$ has exactly $\varphi(p-1)$ primitive roots.
\end{enumerate}
\end{theorem}
\begin{proof}
(a) Write $m = qd + r$ with $0 \le r < d = \ord_n(a)$; then $a^r \equiv a^m \equiv 1 \implies r = 0$.
(b) For prime $p$ and $d \mid p-1$, if an element $a$ of order $d$ exists, then $1, a, \dots, a^{d-1}$ are all $d$ roots of $x^d - 1 \equiv 0 \pmod p$, and the elements of order $d$ are $a^k$ with $\gcd(k,d)=1$, so there are $\psi(d) = \varphi(d)$ of them. Thus $\psi(d) \le \varphi(d)$ for all $d \mid p-1$; summing over $d \mid p-1$ gives $p-1 = \sum \psi(d) \le \sum \varphi(d) = p-1$, forcing $\psi(d) = \varphi(d)$ for all $d$, including $d = p-1$.
\end{proof}

\begin{example}
Find the last two digits of $7^{2026}$.
\end{example}
\begin{solution}
Modulo $100$: $7^2 = 49$, $7^4 = 2401 \equiv 1 \pmod{100}$, so $\ord_{100}(7) = 4$. Since $2026 \equiv 2 \pmod 4$, $7^{2026} \equiv 7^2 = 49$.
\end{solution}

\begin{problem}
Find all positive integers $n$ such that $n \mid 2^n - 1$.
\end{problem}
\begin{solution}
Only $n = 1$. Suppose $n > 1$ works, and let $p$ be the \emph{smallest} prime
divisor of $n$ ($n$ is odd). Then $2^n \equiv 1 \pmod p$, so $d = \ord_p(2)$ divides both $n$ and $p-1$. Since $n$ has no prime factor $< p$, $\gcd(n, p-1) = 1$, forcing $d = 1$ and $p \mid 2^1 - 1 = 1$, contradiction.
\end{solution}

\section{Prime Valuations, LTE, Lucas, and Cyclotomic Polynomials}\label{sec:valuations}

\begin{definition}
For a prime $p$ and nonzero integer $x$, $\nu_p(x)$ is the exponent of $p$ in
the prime factorization of $x$ ($\nu_p(0) = +\infty$).
\end{definition}

\begin{theorem}[Valuations, Legendre's, Kummer's, and Lucas's Theorems]
\leavevmode
\begin{enumerate}[label=(\alph*)]
\item $\nu_p(xy) = \nu_p(x)+\nu_p(y)$, and $\nu_p(x + y) \ge \min\{\nu_p(x), \nu_p(y)\}$ with \emph{equality} whenever $\nu_p(x) \ne \nu_p(y)$.
\item \textbf{(Legendre)} $\displaystyle \nu_p(n!) = \sum_{i \ge 1}\left\lfloor \frac{n}{p^i}\right\rfloor = \frac{n - s_p(n)}{p-1}$, where $s_p(n)$ is the digit sum of $n$ in base $p$.
\item \textbf{(Kummer)} $\nu_p\binom{m+n}{m} = \frac{s_p(m)+s_p(n)-s_p(m+n)}{p-1}$ equals the number of carries when adding $m$ and $n$ in base $p$.
\item \textbf{(Lucas)} For $n = \sum n_i p^i$ and $k = \sum k_i p^i$ in base $p$, $\binom{n}{k} \equiv \prod_i \binom{n_i}{k_i} \pmod p$ (expand $(1+x)^n = \prod (1+x^{p^i})^{n_i}$ in $\FF_p[x]$).
\end{enumerate}
\end{theorem}

\begin{theorem}[Lifting the Exponent (LTE)]
Let $p$ be an odd prime, $p \mid a - b$, and $p \nmid a, b$. Then
\[
\nu_p(a^n - b^n) = \nu_p(a - b) + \nu_p(n) \qquad \bigl(\text{and for odd } n \text{ with } p\mid a+b:\ \nu_p(a^n+b^n) = \nu_p(a+b) + \nu_p(n)\bigr).
\]
For $p = 2$ and odd $a, b$: if $4 \mid a - b$, $\nu_2(a^n - b^n) = \nu_2(a-b) + \nu_2(n)$;
if $n$ is even, $\nu_2(a^n-b^n) = \nu_2(a-b) + \nu_2(a+b) + \nu_2(n) - 1$.
\end{theorem}
\begin{proof}[Proof for odd $p$]
\emph{Step 1 ($p \nmid n$):} $a^n - b^n = (a-b)(a^{n-1} + \dots + b^{n-1})$, and the second factor is $\equiv n a^{n-1} \not\equiv 0 \pmod p$, so $\nu_p(a^n - b^n) = \nu_p(a-b)$.
\emph{Step 2 ($n = p$):} Write $b = a + pk$. Modulo $p^2$,
\[
\sum_{i=0}^{p-1} a^i b^{p-1-i} \equiv p\,a^{p-1} + pk\,a^{p-2}\frac{p(p-1)}{2} \equiv p\,a^{p-1} \pmod{p^2},
\]
so $\nu_p(a^p - b^p) = \nu_p(a-b) + 1$.
\emph{Step 3:} Write $n = p^t m$ ($p \nmid m$), apply Step 2 $t$ times, then Step 1.
\end{proof}

\begin{theorem}[Zsigmondy, Cyclotomic Polynomials, and Wolstenholme]
\leavevmode
\begin{enumerate}[label=(\alph*)]
\item \textbf{(Zsigmondy)} For coprime $a > b \ge 1$ and $n \ge 2$, $a^n - b^n$ has a primitive prime divisor $p \equiv 1 \pmod n$ except when $(a,b,n)=(2,1,6)$ or $n=2, a+b=2^k$ (and $a^n+b^n$ except $2^3+1^3=9$).
\item \textbf{(Cyclotomic Divisors)} Let $\Phi_n(x) = \prod_{\gcd(k,n)=1}(x - e^{2\pi i k/n}) \in \ZZ[x]$, so $x^n - 1 = \prod_{d\mid n}\Phi_d(x)$. Any prime $p \mid \Phi_n(a)$ satisfies either $p \mid n$ or $p \equiv 1 \pmod n$ ($\ord_p(a) = n$). \emph{Proof:} If $d = \ord_p(a) < n$, then $x^n - 1 = (x^d - 1)\Phi_n(x)Q(x)$ has $a$ as a double root mod $p$, so $p \mid n a^{n-1} \implies p \mid n$.
\item \textbf{(Wolstenholme)} For prime $p \ge 5$, the numerator of $\sum_{k=1}^{p-1}\frac{1}{k}$ is divisible by $p^2$ (pair $\frac{1}{k}+\frac{1}{p-k} = \frac{p}{k(p-k)} \equiv -p k^{-2} \pmod{p^2}$), and $\binom{2p-1}{p-1} \equiv 1 \pmod{p^3}$.
\end{enumerate}
\end{theorem}

\begin{problem}
For odd $n \ge 1$, compute $\nu_3(2^n + 1)$, deduce $3^n \nmid 2^n + 1$ for $n \ge 2$, and show every prime $p > 3$ dividing $a^2+a+1$ satisfies $p \equiv 1 \pmod 3$.
\end{problem}
\begin{solution}
By LTE with $p=3, a=2, b=-1$: $\nu_3(2^n - (-1)^n) = \nu_3(2+1) + \nu_3(n) = 1 + \nu_3(n) < n$ for all $n \ge 2$. Also $a^2+a+1 = \Phi_3(a)$, so $p \mid \Phi_3(a) \implies p = 3$ or $p \equiv 1 \pmod 3$.
\end{solution}

\section{Diophantine Equations, Infinite Descent, and Vieta Jumping}

\begin{theorem}[Chicken McNugget, Pythagorean Triples, Pell's Equation, and Hensel]
\leavevmode
\begin{enumerate}[label=(\alph*)]
\item \textbf{(Chicken McNugget / Frobenius)} For coprime $a, b > 0$, the largest integer not of the form $ax + by$ ($x,y \ge 0$) is $ab - a - b$, and $\frac{(a-1)(b-1)}{2}$ nonnegative integers are non-representable.
\item \textbf{(Pythagorean Triples)} Primitive solutions to $a^2+b^2=c^2$ ($2\mid b$) are $a=m^2-n^2, b=2mn, c=m^2+n^2$ for coprime $m>n\ge 1$ of opposite parity.
\item \textbf{(Pell \& Continued Fractions)} For non-square $d > 0$, all positive solutions to $x^2 - dy^2 = 1$ are given by $x_n + y_n\sqrt{d} = (x_1 + y_1\sqrt{d})^n$, and $x_n/y_n$ are convergents of the continued fraction of $\sqrt{d}$.
\item \textbf{(Hensel's Lemma)} If $P(r) \equiv 0 \pmod{p^k}$ and $p \nmid P'(r)$, there is a unique lift $r' \equiv r \pmod{p^k}$ mod $p^{k+1}$ with $P(r') \equiv 0 \pmod{p^{k+1}}$.
\end{enumerate}
\end{theorem}

\begin{problem}[IMO 1988 Problem 6 --- Vieta Jumping]
Let $a, b \in \NN$ be such that $ab + 1 \mid a^2 + b^2$. Prove that $k = \frac{a^2 + b^2}{ab + 1}$ is a perfect square.
\end{problem}
\begin{solution}
Suppose $k$ is not a square. Among all nonnegative integer pairs $(a,b)$ with $\frac{a^2+b^2}{ab+1} = k$, choose one minimizing $a + b$, with $a \ge b \ge 0$. Note $b > 0$ (else $k = a^2$). The quadratic $x^2 - kbx + (b^2 - k) = 0$ has one root $a$; let the other root be $a'$. By Vieta,
\[
a + a' = kb \qquad\text{and}\qquad a a' = b^2 - k.
\]
Thus $a' = kb - a \in \ZZ$, and $a' \ne 0$ since $k \ne b^2$. If $a' \le -1$, then $a'^2 - kba' + b^2 - k \ge 1 + kb + b^2 - k > 0$, impossible. Hence $a' \ge 1$, and $a' = \frac{b^2-k}{a} < \frac{b^2}{a} \le a \implies a' + b < a + b$, contradicting minimality of $a+b$.
\end{solution}

\section{Dirichlet's Theorem, Sums of Squares, and Quadratic Rings}

\begin{theorem}[Dirichlet's Theorem]
If $\gcd(a, n) = 1$, the progression $a, a+n, a+2n, \dots$ contains infinitely many primes.
\end{theorem}

\begin{lemma}[Thue's Lemma]
If $n > 1$ and $\gcd(a, n) = 1$, there exist integers $x, y$ with $0 < |x|, |y| \le \sqrt{n}$ and $ax \equiv y \pmod n$.
\end{lemma}
\begin{proof}
With $m = \lfloor\sqrt{n}\rfloor$, there are $(m+1)^2 > n$ pairs $(x,y) \in \{0,\dots,m\}^2$, so two have $ax_1 - y_1 \equiv ax_2 - y_2 \pmod n$; set $x = x_1-x_2, y = y_1-y_2$.
\end{proof}

\begin{theorem}[Two, Three, and Four Squares, and Gaussian/Eisenstein Integers]
\leavevmode
\begin{enumerate}[label=(\alph*)]
\item \textbf{(Fermat's Two-Square Theorem)} Prime $p = x^2+y^2 \iff p=2$ or $p \equiv 1 \pmod 4$; $n = x^2+y^2 \iff 2 \mid \nu_q(n)$ for every prime $q \equiv 3 \pmod 4$. \emph{Proof for $p \equiv 1 \pmod 4$:} By Wilson, $-1 \equiv \bigl((\frac{p-1}{2})!\bigr)^2 = a^2 \pmod p$. By Thue, $ax \equiv y \pmod p$ with $0 < |x|,|y| < \sqrt{p}$, so $p \mid x^2+y^2 < 2p \implies x^2+y^2=p$.
\item \textbf{(Legendre's Three-Square \& Lagrange's Four-Square)} $n = x^2+y^2+z^2 \iff n \ne 4^a(8b+7)$; every $n \ge 0$ is a sum of four squares.
\item \textbf{(Gaussian \& Eisenstein UFDs)} $\ZZ[i]$ ($N(a+bi)=a^2+b^2$) and $\ZZ[\omega]$ ($\omega = e^{2\pi i/3}, N(a+b\omega)=a^2-ab+b^2$) are Euclidean UFDs.
\end{enumerate}
\end{theorem}

\begin{problem}
(a) Prove that a prime $q \equiv 3 \pmod 4$ dividing $x^2+y^2$ must divide both $x$ and $y$.
(b) Find all integer solutions to $x^2 + 1 = y^3$.
\end{problem}
\begin{solution}
(a) If $q \nmid x$, then $(xy^{-1})^2 \equiv -1 \pmod q \implies 1 \equiv (-1)^{(q-1)/2} = -1 \pmod q$, contradiction.
(b) $x$ must be even and $y$ odd. In $\ZZ[i]$, $(x+i)(x-i) = y^3$; since $\gcd(x+i, x-i) \mid 2i$ and $N(x+i) = y^3$ is odd, $x+i$ and $x-i$ are coprime, so $x+i = (a+bi)^3 \implies 1 = b(3a^2-b^2) \implies b=-1, a=0 \implies (x,y)=(0,1)$.
\end{solution}

\section{Quadratic Residues, Jacobi Symbol, and Reciprocity [!!]}

\begin{definition}
For odd prime $p \nmid a$, the \textbf{Legendre symbol} is $\left(\frac{a}{p}\right) = +1$ if $x^2 \equiv a \pmod p$ is solvable, $-1$ if not. For odd $n = \prod p_i^{a_i}$, the \textbf{Jacobi symbol} is $\left(\frac{a}{n}\right) = \prod \left(\frac{a}{p_i}\right)^{a_i}$.
\end{definition}

\begin{theorem}[Euler's Criterion and Quadratic/Jacobi Reciprocity]
For odd prime $p$, $\left(\frac{a}{p}\right) \equiv a^{(p-1)/2} \pmod p$ and there are $\frac{p-1}{2}$ QRs. For coprime odd positive integers $m, n$:
\[
\left(\frac{m}{n}\right)\left(\frac{n}{m}\right) = (-1)^{\frac{m-1}{2}\cdot\frac{n-1}{2}}, \qquad
\left(\frac{-1}{n}\right) = (-1)^{(n-1)/2}, \qquad
\left(\frac{2}{n}\right) = (-1)^{(n^2-1)/8}.
\]
\end{theorem}

\begin{example}
$\left(\frac{37}{103}\right) = \left(\frac{103}{37}\right) = \left(\frac{29}{37}\right) = \left(\frac{37}{29}\right) = \left(\frac{8}{29}\right) = \left(\frac{2}{29}\right)^3 = (-1)^3 = -1$, so $x^2 \equiv 37 \pmod{103}$ has no solution.
\end{example}

\begin{problem}
Show that for an odd prime $p$, every prime divisor $q$ of $2^p - 1$
satisfies $q \equiv 1 \pmod{2p}$ and $q \equiv \pm 1 \pmod 8$.
\end{problem}
\begin{solution}
$\ord_q(2) = p \mid q-1$ and $2 \mid q-1 \implies 2p \mid q-1$. Also $2 \equiv 2^{p+1} = \bigl(2^{(p+1)/2}\bigr)^2 \pmod q$, so $\left(\frac{2}{q}\right) = 1 \implies q \equiv \pm 1 \pmod 8$.
\end{solution}

\section{Prime Distribution, Approximation, and EGZ [Advanced]}

\begin{theorem}[Prime Number Theorem, Bertrand, Dirichlet Approximation, and EGZ]
\leavevmode
\begin{enumerate}[label=(\alph*)]
\item \textbf{(PNT \& Bertrand)} $\pi(x) \sim \frac{x}{\ln x}$, $\prod_{p \le n} p < 4^n$, and for every $n \ge 1$ there is a prime $p \in (n, 2n]$.
\item \textbf{(Dirichlet's Approximation)} For $\alpha \in \RR$ and $N \ge 1$, there exist $p, q \in \ZZ$ with $1 \le q \le N$ and $|q\alpha - p| < \frac{1}{N}$ (by pigeonhole on $\{0\alpha\}, \{1\alpha\}, \dots, \{N\alpha\} \in [0,1)$).
\item \textbf{(Erdős--Ginzburg--Ziv)} Among any $2n-1$ integers, some $n$ have sum divisible by $n$ (multiplicative in $n$; for prime $p$, sort $0 \le a_1 \le \dots \le a_{2p-1} < p$: either $a_i = a_{i+p-1}$ or Cauchy--Davenport on $\{a_{2p-1}\} + \sum_{i=1}^{p-1}\{a_i, a_{i+p-1}\}$ covers all $p$ residues mod $p$).
\end{enumerate}
\end{theorem}

\begin{problem}
Prove that $H_n = 1 + \frac12 + \dots + \frac1n$ is never an integer for $n \ge 2$.
\end{problem}
\begin{solution}
Let $2^k$ be the largest power of $2$ not exceeding $n$ ($k \ge 1$). Any other multiple of $2^k$ is $\ge 2^{k+1} > n$, so $m = 2^k$ is the \emph{unique} denominator in $\{1,\dots,n\}$ with $\nu_2(m) = k$. Hence in $H_n = \sum_{m=1}^n \frac{1}{m}$, exactly one term has $\nu_2 = -k$ and all others have $\nu_2 > -k$, giving $\nu_2(H_n) = -k < 0 \implies H_n \notin \ZZ$. (Alternatively, by Bertrand's postulate there is a prime $p \in (n/2, n]$, giving $\nu_p(H_n) = -1 < 0$.)
\end{solution}

%======================================================================
\part*{Appendix: How to Train}
\addcontentsline{toc}{part}{Appendix: How to Train}

\begin{itemize}
\item \textbf{Solve slowly, review deeply.} One problem fully understood ---
including \emph{why} the key idea works and where else it applies --- beats
five problems rushed.
\item \textbf{Keep a technique journal.} After each solved (or failed)
problem, record the pivotal idea in one sentence: ``smallest prime factor $+$
order,'' ``tangent line trick,'' ``invert at the tangency point,'' ``Vieta jump on minimal $a+b$.''
\item \textbf{Reconstruct proofs from memory.} Every theorem proved in these
notes should be reproducible on a blank page; the proofs themselves are
exemplars of technique (Cauchy's forward--backward induction, Kelly's extremal proof of Sylvester--Gallai, Thue's
pigeonhole, Fact 5, the Menelaus proof of Pascal).
\item \textbf{Sources of problems.} Past IMO, IMO Shortlists, and national olympiads; work chronologically within a topic, oldest first.
\item \textbf{Mock exams.} Under real conditions: 3 problems, 4.5 hours,
full write-ups. Grading your own write-up against an official solution is
where most improvement happens.
\end{itemize}
\end{document}